%  Методы оптимизации — экзаменационные билеты
%  Компиляция: pdflatex tickets.tex  (дважды, из-за оглавления)
%
%  Это НЕ конспект. Каждый билет самодостаточен и рассчитан на то, чтобы его
%  можно было переписать от руки за 40 минут: ~2 печатные страницы,
%  12–15 выносных формул, без рисунков и численных примеров.
%  Подробная версия — main.tex.

\documentclass[12pt,a4paper]{article}

%  Преамбула конспекта «Методы оптимизации»
%  Компилировать pdfLaTeX (кодировка utf8, шрифты T2A).

% Latin Modern не имеет T2A-начертаний, поэтому кириллица идёт через cm-super:
% семейство cmr в кодировке T2A.
\usepackage[T2A]{fontenc}
\usepackage[utf8]{inputenc}
\usepackage[english,russian]{babel}

\usepackage{amsmath,amssymb,amsthm,mathtools}
\usepackage{geometry}
\geometry{a4paper,margin=2.5cm}

\usepackage{tikz}
\usetikzlibrary{arrows.meta,patterns,decorations.pathmorphing,decorations.pathreplacing,shapes.geometric,calc,positioning}
\usepackage{pgfplots}
\usepgfplotslibrary{groupplots,fillbetween}
\pgfplotsset{compat=1.18}

\usepackage[most]{tcolorbox}
\tcbuselibrary{breakable}

\usepackage{enumitem}
\usepackage{booktabs}
\usepackage{hyperref}
\hypersetup{
  colorlinks=true,
  linkcolor=blue!55!black,
  citecolor=blue!55!black,
  urlcolor=blue!55!black,
  pdftitle={Методы оптимизации --- конспект}
}

%  Окружения для теорем
\theoremstyle{plain}
% Счётчики независимы: теоремы, леммы и утверждения нумеруются каждое своим рядом.
\newtheorem{theorem}{Теорема}
\newtheorem{proposition}{Утверждение}
\newtheorem{corollary}{Следствие}
\newtheorem*{corollary*}{Следствие}
\newtheorem{lemma}{Лемма}
\newtheorem{criterion}{Признак}
\newtheorem*{criterion*}{Признак}

\theoremstyle{definition}
\newtheorem{definition}{Определение}
\newtheorem{example}{Пример}
\newtheorem{exercise}{Упражнение}

\theoremstyle{remark}
\newtheorem*{remark}{Замечание}
\newtheorem*{reminder}{Напоминание}

\renewcommand{\qedsymbol}{$\blacksquare$}

%  Пометки конспекта
%    \begin{idea} ... \end{idea} — пояснения, которых не было на лекции
%    \begin{warn} ... \end{warn} — неразборчиво / описка / требует сверки
\newtcolorbox{idea}[1][]{
  breakable,
  enhanced,
  colback=blue!4!white,
  colframe=blue!45!black,
  boxrule=0.4pt,
  left=6pt,right=6pt,top=5pt,bottom=5pt,
  title={Пояснение},
  fonttitle=\bfseries\small,
  coltitle=white,
  attach boxed title to top left={xshift=6pt,yshift=-2pt},
  boxed title style={colback=blue!45!black,boxrule=0pt,arc=1pt},
  #1
}

\newtcolorbox{warn}[1][]{
  breakable,
  enhanced,
  colback=orange!6!white,
  colframe=orange!65!black,
  boxrule=0.4pt,
  left=6pt,right=6pt,top=5pt,bottom=5pt,
  title={Сверить с записью},
  fonttitle=\bfseries\small,
  coltitle=white,
  attach boxed title to top left={xshift=6pt,yshift=-2pt},
  boxed title style={colback=orange!65!black,boxrule=0pt,arc=1pt},
  #1
}

%  Пошаговая запись метода: \begin{algo}[title={Метод ...}] ... \end{algo}
\newtcolorbox{algo}[1][]{
  breakable,
  enhanced,
  colback=black!3!white,
  colframe=black!60,
  boxrule=0.4pt,
  left=6pt,right=6pt,top=5pt,bottom=5pt,
  title={Алгоритм},
  fonttitle=\bfseries\small,
  coltitle=white,
  attach boxed title to top left={xshift=6pt,yshift=-2pt},
  boxed title style={colback=black!60,boxrule=0pt,arc=1pt},
  #1
}

% Шаги внутри algo
\newlist{steps}{enumerate}{1}
\setlist[steps]{label=\textbf{Шаг \arabic*.}, leftmargin=*, labelsep=0.6em,
                topsep=4pt, itemsep=3pt}

% Отметка начала лекции
\newcommand{\lecture}[2]{%
  \bigskip
  \noindent\hfill\textit{\small #1\ifx\relax#2\relax\else, #2\fi}\\[-0.3em]
  \noindent\rule{\textwidth}{0.4pt}
  \medskip
}

%  Макросы
\newcommand{\R}{\mathbb{R}}
\newcommand{\NN}{\mathbb{N}}
\newcommand{\eps}{\varepsilon}
\newcommand{\T}{^{\mathsf{T}}}
\newcommand{\rank}{\operatorname{rank}}
\newcommand{\sgn}{\operatorname{sign}}
\newcommand{\norm}[1]{\left\lVert #1 \right\rVert}
\newcommand{\abs}[1]{\left\lvert #1 \right\rvert}
\DeclareMathOperator*{\argmin}{arg\,min}
\DeclareMathOperator*{\argmax}{arg\,max}

% Индексные множества канонической задачи.
%   N — индексы переменных, M — индексы ограничений,
%   N_0^+ / N_0^0 — положительные / нулевые компоненты допустимой точки,
%   N^+ — базисное множество (пополненное до |M|).
\newcommand{\Np}{N^{+}}
\newcommand{\Nz}{N^{0}}
\newcommand{\Nzp}{N_0^{+}}
\newcommand{\Nzz}{N_0^{0}}

% Сокращение «опорный вектор», как в тетради.
\newcommand{\ov}{о/в}

% Клетка транспортной таблицы: #1 — тариф c_ij (мелко, сверху),
% #2 — объём перевозки x_ij (обычным кеглем, снизу).
\newcommand{\tc}[2]{\shortstack[r]{\scriptsize $#1$\\ $#2$}}
% Пустая (незаполненная) клетка таблицы.
\newcommand{\tcx}[1]{\tc{#1}{\times}}

% Маркер «++ в кружке» из тетради: доказательство будет дописано позже.
\newcommand{\plusplus}{%
  \tikz[baseline=(pp.base)]{%
    \node[draw, ellipse, inner xsep=3pt, inner ysep=1pt, line width=0.4pt] (pp)
      {\footnotesize $+\,+$};}%
}

%  Быстрая сборка: \showfiguresfalse выключает все TikZ-рисунки.
\newif\ifshowfigures
\showfigurestrue

\newcommand{\inputfig}[1]{%
  \ifshowfigures
    \input{#1}%
  \else
    \begin{center}\fbox{\small рисунок выключен: \texttt{\detokenize{#1}}}\end{center}%
  \fi
}


% Заголовки билетов не нумеруются автоматически: номер стоит в самом названии.
\setcounter{secnumdepth}{0}

% Формулировка вопроса под заголовком билета.
\newcommand{\ticketq}[1]{%
  \noindent\textit{\small #1}\par
  \smallskip
  \noindent\rule{\textwidth}{0.4pt}
  \medskip
}

% Отсылка к подробному разбору в конспекте.
\newcommand{\seefull}[1]{%
  \par\medskip
  \noindent\textit{\footnotesize Подробно: конспект, #1.}%
}

% Текст плотный и насыщен формулами в строке: даём TeX дополнительный запас
% на растяжку, иначе строки переполняются.
\emergencystretch=3em

% Плотный список шагов алгоритма: без вертикальных зазоров.
\newlist{csteps}{enumerate}{1}
\setlist[csteps]{label=\textbf{\arabic*.}, leftmargin=1.5em, labelsep=0.4em,
                 topsep=2pt, itemsep=0pt, parsep=0pt, partopsep=0pt}

% Плотные заголовки-абзацы.
\makeatletter
\renewcommand\paragraph{\@startsection{paragraph}{4}{0pt}%
  {1.2ex \@plus 0.3ex \@minus 0.2ex}%
  {-0.6em}%
  {\normalfont\normalsize\bfseries}}
\makeatother

\title{\textbf{Методы оптимизации}\\[0.3em]\large Ответы на экзаменационные билеты}
\author{}
\date{}

\begin{document}

\maketitle
\thispagestyle{empty}

\begin{idea}[title={Как этим пользоваться}]
Каждый билет — готовый письменный ответ: постановка, определения, формулировки
с доказательствами, алгоритм. Доказательства сжаты до несущей цепочки —
выкладки, без которых доказательство перестаёт быть доказательством,
сохранены; рутинные раскрытия свёрнуты.

\medskip
Численных примеров, рисунков и пояснений здесь нет — они остались в конспекте
(\texttt{main.pdf}), на который каждый билет ссылается в конце.
\end{idea}

\tableofcontents
\newpage

\section{Билет 1. Задача ЛП в канонической форме}

\ticketq{Задача ЛП в канонической форме. Основные понятия. Лемма о числе
опорных векторов. Теорема о совпадении крайней точки и опорного вектора.
Теорема об оптимальном решении задачи ЛП. Метод перебора крайних точек.}

\paragraph{Постановка.}
\[
  c\T[N]x[N] \to \min , \qquad A[M,N]x[N] = b[M] , \qquad x[N] \ge 0 ,
\]
$\abs{M} = m$, $\abs{N} = n$, $m < n$, $\rank A = m$;
$S = \{\, x[N] \mid A[M,N]x[N] = b[M],\ x[N] \ge 0 \,\}$ — многогранное
допустимое множество.

\paragraph{Основные понятия.}
$x_0 \in S$ — \textbf{крайняя точка}, если
$x_0 \ne (1-\lambda)x_1 + \lambda x_2$ ни при каких $x_1 \ne x_2 \in S$,
$0 < \lambda < 1$.

$\Nzp = \{j : x_0[j] > 0\}$, $\Nzz = \{j : x_0[j] = 0\}$
$\Rightarrow$ $\sum_{j \in \Nzp} x_0[j]\,a_j[M] = b[M]$ ($a_j$ — столбцы $A$).

$x_0$ — \textbf{опорный}, если $\{a_j[M]\}_{j \in \Nzp}$ линейно независимы;
они в $\R^m$ $\Rightarrow$ $\abs{\Nzp} \le m$. Невырожденный:
$\abs{\Nzp} = m$. Вырожденный: $\abs{\Nzp} < m$, тогда $\Nzp \subset N_0$,
$\abs{N_0} = m$, $\det A[M,N_0] \ne 0$, $x_0[N_0 \setminus \Nzp] = 0$.
\textbf{Базис} — столбцы $a_j$, $j \in N_0$.

\paragraph{Лемма (о числе опорных векторов).}
\textit{Опорных векторов не больше $\binom{n}{m}$.}

\textsc{Док-во.} Опорный вектор $\mapsto$ базис $=$ набор $m$ столбцов из $n$;
таких наборов $\binom{n}{m}$. Для набора $N_k$ с $\det A[M,N_k] \ne 0$ решение
системы $A[M,N_k]x[N_k] = b[M]$ единственно и дополняется нулями
$\Rightarrow$ разных опорных векторов не больше числа наборов.
$\blacksquare$

\paragraph{Теорема (крайняя точка $=$ опорный вектор).}
\textit{$x[N] \in S$ опорный $\iff$ $x[N]$ — крайняя точка $S$.}
Далее $\Np = \{j : x[j] > 0\}$.

\textsc{Док-во. 1) Не опорный $\Rightarrow$ не крайняя.}
\begin{gather*}
  \exists\, u[\Np] \ne 0 :\ A[M,\Np]u[\Np] = 0 , \qquad
  u[N \setminus \Np] := 0 \\
  \Rightarrow\quad A[M,N]u[N] = 0
  \quad\Rightarrow\quad A(x \pm \theta u) = b .
\end{gather*}
Неотрицательность держится при $\theta \le \theta_0$:
\[
  \theta_0 = \min\{\theta_1, \theta_2\} , \quad
  \theta_1 = \min_{i \in \Np,\, u[i]>0} \frac{x[i]}{u[i]} , \quad
  \theta_2 = \min_{i \in \Np,\, u[i]<0} \left(-\frac{x[i]}{u[i]}\right) .
\]
$x_{1,2} = x \pm \theta_0 u \in S$, $x_1 \ne x_2$ (ибо $u \ne 0$),
$x = \tfrac12(x_1+x_2)$ $\Rightarrow$ $x$ не крайняя.

\textsc{2) Не крайняя $\Rightarrow$ не опорный.}
$x = (1-\lambda)x_1 + \lambda x_2$. Для $i \in \Nz$:
\[
  0 = \underbrace{(1-\lambda)x_1[i]}_{\ge\,0}
    + \underbrace{\lambda x_2[i]}_{\ge\,0}
  \ \Rightarrow\ x_1[i] = x_2[i] = 0
\]
$\Rightarrow$ $x_1, x_2$ различаются лишь на $\Np$. Из $Ax_1 = Ax_2 = b$:
\[
  A[M,\Np]\big(x_1[\Np] - x_2[\Np]\big) = 0 ,
  \quad x_1[\Np] - x_2[\Np] \ne 0
  \ \Rightarrow\ \text{столбцы зависимы.} \qquad \blacksquare
\]

\paragraph{Теорема (об оптимальном решении).}
\textit{Если задача ЛП разрешима, то хотя бы одно оптимальное решение —
опорный вектор.}

\textsc{Док-во.} Берём оптимальное $\bar{x}$ с минимальным
$\abs{\overline{N}{}^+}$, где $\overline{N}{}^+ = \{j : \bar{x}[j] > 0\}$.
Пусть $\bar{x}$ не опорный:
\[
  \exists\, u[\overline{N}{}^+] \ne 0 :\ A[M,\overline{N}{}^+]u = 0 ,
  \quad u[N \setminus \overline{N}{}^+] := 0
  \ \Rightarrow\ A[M,N]u[N] = 0 .
\]
Тогда $x_\theta := \bar{x} + \theta u$ даёт $Ax_\theta = b$ и $x_\theta \ge 0$
при $\abs{\theta} \le \theta_0$, то есть $x_\theta \in S$, причём
\[
  c\T x_\theta = c\T\bar{x} + \theta\, c\T u
  \ \Rightarrow\
  \big[\, c\T u \ne 0 \ \Rightarrow\ \exists\,\theta :
    c\T x_\theta < c\T\bar{x} \,\big]
  \ —\ \text{противоречие.}
\]
Значит $c\T u = 0$ и все $x_\theta$ оптимальны. При $\theta = \theta_1$
(либо $-\theta_2$) одна положительная компонента обнуляется, новых не
возникает ($u = 0$ вне $\overline{N}{}^+$)
$\Rightarrow$ оптимум с меньшим числом положительных компонент —
противоречие с выбором $\bar{x}$. $\blacksquare$

\textbf{Следствие.} Оптимум достигается в крайней точке; их конечное число.

\paragraph{Метод перебора крайних точек.}
\begin{steps}
  \item Взять $N_k \subset N$, $\abs{N_k} = m$.
  \item $\det A[M,N_k] = 0$ — отбросить.
  \item Решить $A[M,N_k]x[N_k] = b[M]$, положить $x[N \setminus N_k] = 0$.
  \item $x[N_k] \ge 0$ — крайняя точка, считаем $F = c\T[N]x[N]$; иначе
        отбросить.
  \item Перебрать все наборы, взять минимальное $F$.
\end{steps}

Конечен по лемме; нет допустимых точек $\Rightarrow$ $S = \varnothing$.
Непригоден практически: $\binom{50}{25} \approx 1{,}3\cdot10^{14}$, и
неограниченность не распознаётся $\Rightarrow$ симплекс-метод.

\seefull{гл.~1, \S1--\S6 и \S8}

\section{Билет 2. Симплекс-метод. Искусственный базис}

\ticketq{Симплекс-метод. Метод искусственного базиса.}

\paragraph{Обозначения.}
Каноническая задача $c\T[N]x[N] \to \min$, $Ax = b$, $x \ge 0$. Метод строит
опорные векторы $x^0, x^1, \dots$ с невозрастающим $c\T x^k$ (оптимум
достигается в опорном векторе — билет~1).

На $k$-й итерации: базис $N_k$, $\abs{N_k} = m$, $\det A[M,N_k] \ne 0$;
$B[N_k,M] = A[M,N_k]^{-1}$; $L_k = N \setminus N_k$, $x^k[L_k] = 0$.
Потенциалы и оценки:
\[
  A\T[M,N_k]\,y^k[M] = c[N_k] ,
  \qquad
  d[j] = c[j] - (y^k)\T a_j[M] , \quad j \in L_k .
\]

\paragraph{Теорема (критерий оптимальности).}
\textit{$d[j] \ge 0$ $\forall j \in L_k$ $\Rightarrow$ $x^k$ оптимален.}

\textsc{Док-во.} $x \in S$. Из $Ax = Ax^k = b$, разбивая на базисную и
небазисную части:
\begin{gather*}
  A[M,N_k]\big(x[N_k] - x^k[N_k]\big) + A[M,L_k]x[L_k] = 0 \\
  \Rightarrow\quad x[N_k] - x^k[N_k] = -B[N_k,M]A[M,L_k]x[L_k] .
\end{gather*}
Подставляя и учитывая $(y^k)\T = c\T[N_k]B[N_k,M]$:
\[
  c\T x - c\T x^k
   = c\T[N_k]\big(x[N_k] - x^k[N_k]\big) + c\T[L_k]x[L_k]
   = \sum_{j \in L_k} d[j]\,x[j] \ \ge\ 0 . \qquad \blacksquare
\]

\paragraph{Направление перехода.}
Пусть $\exists\, j_k : d[j_k] < 0$. Берём
\[
  u^k[N_k] = B[N_k,M]a_{j_k}[M] , \quad u^k[j_k] = -1 , \quad
  u^k[L_k \setminus \{j_k\}] = 0 , \quad x_\theta = x^k - \theta u^k .
\]
\textbf{Равенства:} $Au^k = A[M,N_k]B[N_k,M]a_{j_k} - a_{j_k} = 0$ (так как
$A[M,N_k]B = E$) $\Rightarrow$ $Ax_\theta = b$ $\forall\theta$.

\textbf{Цель:} $c\T u^k = (y^k)\T a_{j_k} - c[j_k] = -d[j_k]$, откуда
\[
  F(x_\theta) = F(x^k) + \theta\,d[j_k] < F(x^k) \quad (\theta > 0) .
\]

\paragraph{Выбор шага.}
Требуется $x^k[i] - \theta u^k[i] \ge 0$, $i \in N_k$:
\[
  u^k[i] \le 0 \ \forall i \in N_k \ \Rightarrow\ F(x_\theta) \to -\infty
  \ \ (\text{не ограничена снизу});
\]
иначе $\theta_k = \min\limits_{i \in N_k,\ u^k[i]>0} x^k[i]/u^k[i]$ на
индексе $i_k$; тогда $x_{\theta_k}[i_k] = 0$: $x[i_k]$ выходит, $x[j_k]$ входит.

\paragraph{Алгоритм.}
Дан $x^0$ и базис $N_0$.
\begin{csteps}
  \item $B[N_k,M] = A[M,N_k]^{-1}$; найти $y^k$ из $A\T[M,N_k]y^k = c[N_k]$.
  \item $d[j] = c[j] - (y^k)\T a_j$, $j \in L_k$.
  \item Все $d[j] \ge 0$ — $x^k$ оптимален, конец.
  \item Иначе выбрать $j_k$ с $d[j_k] < 0$, построить $u^k$.
  \item $u^k[i] \le 0$ $\forall i \in N_k$ — не ограничена снизу, конец.
  \item Найти $\theta_k$, $i_k$; $x^{k+1} = x^k - \theta_k u^k$,
        $N_{k+1} = (N_k \setminus \{i_k\}) \cup \{j_k\}$, к~шагу~1.
\end{csteps}

\paragraph{Конечность.}
В невырожденном случае $\theta_k > 0$, и по формуле изменения цели
$F$ \emph{строго} убывает $\Rightarrow$ базисы не повторяются; базисов не
больше $\binom{n}{m}$ (лемма билета~1) $\Rightarrow$ алгоритм конечен.

\paragraph{Вырожденность и зацикливание.}
$\abs{\Nzp} < m$ $\Rightarrow$ возможно $\theta_k = 0$: вектор не меняется,
меняется базис, $F$ не убывает $\Rightarrow$ \textbf{зацикливание} (возврат к
прежнему базису). Средство — \textbf{правило Бленда}: при нескольких
кандидатах на ввод и на вывод брать наименьший индекс; тогда ни один базис не
повторяется и конечность восстанавливается.

\paragraph{Метод искусственного базиса.}
Нужен, когда $x^0$ неизвестен. Добиваемся $b[M] \ge 0$ (строку с $b[i] < 0$
умножаем на $-1$) и решаем
\[
  \sum_{i \in M} y[i] \to \min , \quad
  A[M,N]x[N] + E[M,M]y[M] = b[M] , \quad x \ge 0,\ y \ge 0 ,
\]
старт $x[N] = 0$, $y[M] = b[M]$ — очевидный опорный вектор с базисом $E$.

\paragraph{Теорема (об искусственном базисе).}
\textit{Пусть $(x^*, y^*)$ — оптимум вспомогательной задачи. Тогда
$\sum_i y^*[i] > 0$ $\Rightarrow$ исходная задача несовместна;
$\sum_i y^*[i] = 0$ $\Rightarrow$ $x^*[N] \in S$.}

\textsc{Док-во.} $y \ge 0 \Rightarrow \sum_i y[i] \ge 0$: минимум неотрицателен.

\emph{Случай $>0$}, от противного: $\exists\,\hat{x} \in S$ $\Rightarrow$ пара
$(\hat{x}, 0)$ допустима со значением $0$ — противоречие.

\emph{Случай $=0$:} $y^*[i] \ge 0$ и $\sum_i y^*[i] = 0$ $\Rightarrow$
$y^*[M] = 0$ $\Rightarrow$ $A[M,N]x^*[N] = b[M]$, а $x^* \ge 0$ $\Rightarrow$
$x^* \in S$. $\blacksquare$

\paragraph{Алгоритм искусственного базиса.}
\begin{csteps}
  \item Привести к $b[M] \ge 0$, ввести $y[M] \ge 0$, составить вспомогательную
        задачу.
  \item Решить её симплекс-методом от базиса $x = 0$, $y = b$.
  \item $\sum_i y^*[i] > 0$ — несовместна, конец.
  \item Иначе $y^* = 0$, $x^*$ допустим.
  \item Искусственные переменные, оставшиеся в базисе с нулём, заменить
        столбцами $A[M,N]$ с сохранением линейной независимости; если замены
        нет — уравнение зависимо, вычеркнуть.
  \item Полученный базис — начальный для симплекс-метода на исходной задаче.
\end{csteps}

\seefull{гл.~1, \S7, \S9, \S10 (численные примеры там же)}

\section{Билет 3. Метод потенциалов решения транспортных задач}

\ticketq{Метод потенциалов решения транспортных задач.}

\paragraph{Постановка.}
$m$ поставщиков с запасами $a_i$, $n$ потребителей с потребностями $b_j$,
тарифы $c_{ij}$, перевозки $x_{ij}$:
\[
  F(X) = \sum_{i=1}^{m}\sum_{j=1}^{n} c_{ij}x_{ij} \to \min , \quad
  \sum_{j=1}^{n} x_{ij} = a_i , \quad
  \sum_{i=1}^{m} x_{ij} = b_j , \quad x_{ij} \ge 0 .
\]
Задача \textbf{закрытая}, если $\sum_i a_i = \sum_j b_j$; иначе открытая и
приводится к закрытой фиктивным поставщиком ($a_{m+1} = \sum_j b_j - \sum_i a_i$)
или потребителем ($b_{n+1} = \sum_i a_i - \sum_j b_j$) с нулевыми тарифами.

\paragraph{Баланс необходим и достаточен.}
\textit{Задача разрешима $\iff$ $\sum_i a_i = \sum_j b_j$.}

\textsc{Док-во.} \emph{Необходимость:} суммируя ограничения по строкам и по
столбцам, получаем одну и ту же величину:
$\sum_i a_i = \sum_i\sum_j x_{ij} = \sum_j b_j$.
\emph{Достаточность:} при балансе план существует явно —
$x_{ij} = a_i b_j / \sum_k a_k$ удовлетворяет всем ограничениям.
$\blacksquare$

\paragraph{Базисный план и потенциалы.}
Уравнений $m+n$, одно зависимо (та же сумма по строкам и по столбцам)
$\Rightarrow$ $\rank = m+n-1$ $\Rightarrow$ невырожденный базисный план имеет
ровно $m+n-1$ заполненных клеток. Если клеток меньше, план вырожден: добавляют
базисные клетки с нулевой перевозкой, \emph{так чтобы не образовался цикл}
(иначе набор клеток перестанет быть базисом).

Строке $i$ отвечает потенциал $u_i$, столбцу $j$ — $v_j$; для заполненных
клеток $v_j - u_i = c_{ij}$. Уравнений $m+n-1$, неизвестных $m+n$
$\Rightarrow$ один потенциал задаётся произвольно ($u_1 = 0$), остальные
однозначны. Для пустых клеток $\Delta_{ij} = v_j - u_i - c_{ij}$; условие
оптимальности — $\Delta_{ij} \le 0$.

\paragraph{Теорема (об оптимальности плана).}
\textit{Если для допустимого базисного плана $X$ найдены $u_i, v_j$ с}
\[
  v_j - u_i = c_{ij} \ \text{(заполненные)} , \qquad
  v_j - u_i \le c_{ij} \ \text{(пустые)} ,
\]
\textit{то $X$ оптимален.}

\textsc{Док-во.} Пусть $Y = (y_{ij})$ — любой допустимый план. Для всех клеток
$c_{ij} \ge v_j - u_i$; умножая на $y_{ij} \ge 0$ и суммируя:
\[
  F(Y) \ \ge\ \sum_{i}\sum_{j}(v_j - u_i)y_{ij}
       = \sum_{j} v_j \underbrace{\sum_{i} y_{ij}}_{=\,b_j}
       - \sum_{i} u_i \underbrace{\sum_{j} y_{ij}}_{=\,a_i}
       = \sum_j v_j b_j - \sum_i u_i a_i .
\]
Для самого $X$: на заполненных клетках $c_{ij} = v_j - u_i$, на пустых
$x_{ij} = 0$ $\Rightarrow$ замена $c_{ij} \to v_j - u_i$ законна во всех
клетках, и
\[
  F(X) = \sum_i\sum_j (v_j - u_i)x_{ij} = \sum_j v_j b_j - \sum_i u_i a_i
  \ \Rightarrow\ F(Y) \ge F(X) . \qquad \blacksquare
\]

Величина $\sum_j v_j b_j - \sum_i u_i a_i$ — значение двойственной задачи:
неравенство есть слабая двойственность, совпадение на $X$ — сильная, а проверка
$\Delta_{ij} \le 0$ есть проверка двойственной допустимости.

\paragraph{Цикл пересчёта.}
Если $\exists\,(p,q)$ пустая с $\Delta_{pq} > 0$, она вводится в базис. Цикл
пересчёта — замкнутая ломаная: начинается в $(p,q)$, идёт по строкам и
столбцам, поворачивает только в заполненных клетках, возвращается в $(p,q)$.
Знаки чередуются от $\oplus$ в $(p,q)$: $C_+$ и $C_-$ — клетки со знаками
$\oplus$ и $\ominus$. Величина пересчёта и новый план:
\[
  \theta = \min_{(i,j) \in C_-} x_{ij} ,
  \qquad
  x_{ij} := x_{ij} + \theta \ \ \text{на } C_+ , \qquad
  x_{ij} := x_{ij} - \theta \ \ \text{на } C_- .
\]
Клетка из $C_-$, на которой достигнут минимум, выходит из базиса.

\paragraph{Теорема (об улучшении плана).}
\textit{$\Delta F = -\theta\,\Delta_{pq} < 0$.}

\textsc{Док-во.}
$\Delta F = \theta\big(\sum_{C_+} c_{ij} - \sum_{C_-} c_{ij}\big)$.
Все клетки цикла, кроме $(p,q)$, заполнены $\Rightarrow$ для них
$c_{ij} = v_j - u_i$; для стартовой $c_{pq} = v_q - u_p - \Delta_{pq}$. Цикл
замкнут: каждая задействованная строка и каждый столбец встречаются ровно
дважды — по разу со знаками $\oplus$ и $\ominus$ $\Rightarrow$ все $u_i, v_j$
взаимно сокращаются, остаётся добавка стартовой клетки:
\[
  \sum_{C_+} c_{ij} - \sum_{C_-} c_{ij} = -\Delta_{pq}
  \ \Rightarrow\ \Delta F = -\theta\,\Delta_{pq} < 0
  \quad (\theta > 0,\ \Delta_{pq} > 0) . \qquad \blacksquare
\]
Отсюда правило выбора: брать клетку с \emph{наибольшим} $\Delta_{ij} > 0$.

\paragraph{Алгоритм метода потенциалов.}
\begin{csteps}
  \item Проверить закрытость, при необходимости ввести фиктивного поставщика
        или потребителя.
  \item Построить начальный базисный план (северо-западный угол либо
        минимальный элемент); проверить число клеток $m+n-1$.
  \item Найти потенциалы: $u_1 = 0$, далее $v_j - u_i = c_{ij}$ по заполненным
        клеткам.
  \item Для пустых клеток вычислить $\Delta_{ij} = v_j - u_i - c_{ij}$.
  \item Все $\Delta_{ij} \le 0$ — план оптимален, конец.
  \item Иначе взять клетку с максимальным $\Delta_{ij} > 0$, построить цикл,
        найти $\theta$, пересчитать план, вывести клетку с нулём на минусе;
        к~шагу~3.
\end{csteps}

Метод потенциалов — симплекс-метод в табличной форме: $u_i, v_j$ играют роль
потенциалов $y^k$, $\Delta_{ij}$ — роль оценок $-d[j]$.

\seefull{гл.~2, \S4, \S5, \S7 (численный пример в \S4)}

\section{Билет 4. Начальный план и усложнения ТЗ}

\ticketq{Методы построения начального плана. Усложнения транспортной задачи.}

\paragraph{Что требуется от начального плана.}
Для старта метода потенциалов нужен допустимый \emph{базисный} план закрытой
задачи: $\sum_j x_{ij} = a_i$, $\sum_i x_{ij} = b_j$, $x_{ij} \ge 0$ и ровно
$m+n-1$ заполненных клеток (ранг системы $m+n-1$: одно из $m+n$ уравнений
зависимо). Клеток меньше — план вырожден, добавляют нулевые базисные перевозки.

\paragraph{Метод северо-западного угла.}
Заполнение с клетки $(1,1)$, тарифы не учитываются.
\begin{csteps}
  \item В текущую клетку $(i,j)$ записать $x_{ij} = \min\{a_i, b_j\}$.
  \item $a_i := a_i - x_{ij}$, $b_j := b_j - x_{ij}$.
  \item $a_i = 0$ — к следующей строке; $b_j = 0$ — к следующему столбцу; при
        $a_i = b_j = 0$ закрыть одну строку и один столбец, добавив при
        необходимости нулевую базисную перевозку.
  \item Повторять до распределения всего; проверить число клеток $m+n-1$.
\end{csteps}

\paragraph{Метод минимального элемента.}
То же, но клетка выбирается по тарифу: среди незакрытых берётся $(i,j)$ с
минимальным $c_{ij}$ (при равенстве — самая северо-западная), далее шаги 1--4
дословно те же. Даёт обычно более дешёвый план, так как сначала используются
дешёвые маршруты, но оптимальности не гарантирует.

\paragraph{Оба метода дают базисный план.}
Выкладка общая, различается только правило выбора клетки.
\[
  a_i, b_j \ge 0 \ \Rightarrow\ x_{ij} = \min\{a_i, b_j\} \ge 0 ;
\]
каждый шаг закрывает хотя бы строку или столбец $\Rightarrow$ процесс конечен,
а заполненных клеток не более $m+n-1$ (на последнем шаге закрываются обе);
задача закрытая ($\sum_i a_i = \sum_j b_j$) $\Rightarrow$ исчерпание всех
запасов означает удовлетворение всех потребностей $\Rightarrow$ выполнены
$\sum_j x_{ij} = a_i$, $\sum_i x_{ij} = b_j$, $x_{ij} \ge 0$ — план допустим.

Он к тому же \emph{базисный}: заполненные клетки образуют «ступенчатую»
цепочку (каждый шаг сдвигается на строку вниз или на столбец вправо и назад
не возвращается), поэтому цикла среди них нет, то есть соответствующие столбцы
линейно независимы. $\blacksquare$

\paragraph{Проверка плана.}
$\sum_j x_{ij} = a_i$; $\sum_i x_{ij} = b_j$; $x_{ij} \ge 0$; число базисных
клеток $m+n-1$.

\paragraph{Усложнения транспортной задачи.}
\begin{csteps}
  \item \textbf{Открытая модель} $\sum_i a_i \ne \sum_j b_j$ — приводится к
        закрытой фиктивным участником с нулевыми тарифами:
        \[
          \sum_i a_i > \sum_j b_j :\quad
            b_{n+1} = \sum_i a_i - \sum_j b_j ,\ c_{i,n+1} = 0 ;
        \]
        \[
          \sum_i a_i < \sum_j b_j :\quad
            a_{m+1} = \sum_j b_j - \sum_i a_i ,\ c_{m+1,j} = 0 .
        \]
        Перевозки в фиктивный пункт — это нераспределённый запас (или
        неудовлетворённая потребность); нулевые тарифы не меняют $F$, поэтому
        оптимум исходной задачи сохраняется.
  \item \textbf{Вырожденность:} заполненных клеток меньше $m+n-1$ — добавляют
        нулевые базисные перевозки так, чтобы не образовался цикл.
  \item \textbf{Многопродуктовая модель.} Для продукта $k$ вводятся
        $x_{ij}^{(k)}$:
        $\sum_j x_{ij}^{(k)} = a_i^{(k)}$,
        $\sum_i x_{ij}^{(k)} = b_j^{(k)}$,
        $\sum_k x_{ij}^{(k)} \le d_{ij}$ (пропускная способность маршрута).
  \item \textbf{Запрет поставки} из $A_i$ в $B_j$: клетку исключают либо
        задают $c_{ij} = M$ (большое число). При достаточно большом $M$
        положительная перевозка по этой клетке увеличила бы $F$ сильнее любого
        допустимого варианта, поэтому в оптимальном плане она нулевая — если
        допустимый план без неё вообще существует.
  \item \textbf{Обязательная поставка:} $x_{ij} > 0$, на практике заменяется
        на $x_{ij} \ge K_{ij}$, $K_{ij} > 0$.
  \item \textbf{Не менее $K_{ij}$:} объём выделяют заранее,
        $a_i := a_i - K_{ij}$, $b_j := b_j - K_{ij}$ (сдвиг
        $x_{ij}' = x_{ij} - K_{ij}$) — задача снова обычная.
  \item \textbf{Не более $K_{ij}$:} $x_{ij} \le K_{ij}$ — пропускная
        способность.
  \item \textbf{Сменно-суточный характер:} индекс времени $x_{ij}^{(t)}$;
        учитываются смены, сутки, время доставки, доступность транспорта.
  \item \textbf{Возврат тары, порожний пробег:} обратные перевозки
        $x_{ji}^{\text{back}}$, возможна задача минимизации порожнего пробега.
\end{csteps}

Пункты 1--5 не выводят задачу из класса транспортных: запрет и верхняя граница
— ограничения на отдельные клетки, нижняя снимается сдвигом переменной.
Пункты 6--7 увеличивают размерность и требуют уже общего симплекс-метода, а не
табличной схемы.

\seefull{гл.~2, \S3 и \S8 (численные примеры в \S3)}

\section{Билет 5. Унимодальные функции. Одномерная минимизация без производной}

\ticketq{Унимодальные функции. Методы одномерной минимизации, не использующие
вычисление производной функции цели.}

\paragraph{Постановка.}
$f(x) \to \min$, $x \in [a,b]$. Методы \textbf{нулевого порядка} используют
только значения $f(x_1), f(x_2), \dots$ и не используют $f'(x)$. Применяются
самостоятельно либо как подбор шага $\alpha_k$ в многомерных методах:
$\min_\alpha f(x_k + \alpha d_k)$.

\paragraph{Унимодальность.}
$f$ \textbf{унимодальна} на $[a,b]$, если существует единственная точка
минимума $x_*$, причём $f$ убывает на $[a, x_*]$ и возрастает на $[x_*, b]$:
\[
  x_1 < x_2 \le x_* \Rightarrow f(x_1) \ge f(x_2) ,
  \qquad
  x_* \le x_1 < x_2 \Rightarrow f(x_1) \le f(x_2) .
\]
\textbf{Интервал неопределённости} — отрезок $[a_k, b_k] \ni x_*$,
$L_k = b_k - a_k$. Останов при $L_k \le \delta$, тогда
$x_* \approx (a_k + b_k)/2$.

\paragraph{Лемма (об исключении части отрезка).}
\textit{Пусть $f$ унимодальна на $[a,b]$ и $x_1 < x_2$. Тогда}
\[
  f(x_1) \ge f(x_2) \ \Rightarrow\ x_* \notin [a, x_1) ,
  \qquad
  f(x_1) \le f(x_2) \ \Rightarrow\ x_* \notin (x_2, b] .
\]

\textsc{Док-во.} Первое, от противного: $x_* < x_1 < x_2$ $\Rightarrow$ обе
точки правее минимума $\Rightarrow$ по возрастанию справа
$f(x_1) \le f(x_2)$, а по условию $f(x_1) \ge f(x_2)$ $\Rightarrow$
$f(x_1) = f(x_2)$ — невозможно при единственном минимуме на строго
возрастающем участке. Второе симметрично: $x_1 < x_2 < x_*$ $\Rightarrow$
$f(x_1) \ge f(x_2)$ по убыванию слева, что противоречит
$f(x_1) \le f(x_2)$. $\blacksquare$

Вся дальнейшая работа — применение леммы: по двум сравнённым значениям
отбрасывается часть отрезка.

\paragraph{Равномерный поиск.}
$h = (b-a)/n$, $x_i = a + ih$, $i = \overline{0,n}$. Найти $j$ с
$f(x_j) = \min_i f(x_i)$ $\Rightarrow$ $x_* \in [x_{j-1}, x_{j+1}]$, длина
$2h$. Стоит $n+1$ вычислений.

\paragraph{Локализация экстремума.}
То же, но без просмотра всей сетки: идём слева и останавливаемся на первом
$i$ с
\[
  f\big(a+(i-1)h\big) > f\big(a+ih\big) \le f\big(a+(i+1)h\big)
  \ \Rightarrow\ x_* \in \big[a+(i-1)h,\ a+(i+1)h\big] .
\]
Результат тот же, вычислений в среднем меньше: обход прекращается сразу
после минимума.

\textbf{Вариант с удвоением шага (метод Свенна)} — применяется, когда отрезок
заранее не задан. Выбираем $x_0$, $h > 0$; по значениям
$f(x_0-h), f(x_0), f(x_0+h)$ определяем направление убывания и идём в нём с
удвоением шага:
\[
  x_{k+1} = x_k + 2^k h .
\]
После первого возрастания минимум локализован: $x_* \in [x_{k-1}, x_{k+1}]$.
Выигрыш в том, что отрезок находится за $O(\log)$ шагов, а не перебором сетки.

\paragraph{Дихотомия.}
$x_{1,2} = \dfrac{a_k+b_k}{2} \mp \eps$, $\eps > 0$ мало. По лемме:
\[
  f(x_1) \ge f(x_2) \Rightarrow [a_{k+1},b_{k+1}] = [x_1, b_k] ,
  \qquad
  \text{иначе } [a_k, x_2] .
\]
$L_{k+1} = L_k/2 + \eps$ за \emph{два} вычисления $\Rightarrow$ коэффициент
$\sqrt{1/2} \approx 0{,}707$ на одно вычисление. Параметр $\eps$ берётся на
пределе различимости значений $f$.

\paragraph{Бисекция (деление отрезка пополам).}
Три точки $x_1 = a_k + L_k/4$, $x_2 = (a_k+b_k)/2$, $x_3 = a_k + 3L_k/4$:
\[
  f(x_1) < f(x_2) \Rightarrow [a_k, x_2] ; \quad
  \text{иначе } f(x_2) > f(x_3) \Rightarrow [x_2, b_k] ; \quad
  f(x_2) \le f(x_3) \Rightarrow [x_1, x_3] .
\]
Отрезок ровно вдвое короче, а серединой нового становится одна из уже
вычисленных точек $\Rightarrow$ два новых вычисления на итерацию, тот же
коэффициент $0{,}707$.

\paragraph{Золотое сечение.}
$\lambda_k = a_k + \alpha L_k$, $\mu_k = b_k - \alpha L_k$; правило то же:
$f(\lambda_k) \ge f(\mu_k) \Rightarrow [\lambda_k, b_k]$, иначе
$[a_k, \mu_k]$.

\textsc{Вывод $\alpha$.} После сокращения $L_{k+1} = (1-\alpha)L_k$. Чтобы
одна из старых точек стала внутренней точкой нового отрезка (и её значение
переиспользовалось), нужно
\[
  \frac{\alpha}{1-\alpha} = 1 - \alpha
  \ \Rightarrow\ \alpha = (1-\alpha)^2
  \ \Rightarrow\ \alpha^2 - 3\alpha + 1 = 0
  \ \Rightarrow\ \alpha = \frac{3 - \sqrt5}{2} \approx 0{,}382
\]
(второй корень $>1$ не годится). Итог: \emph{одно} новое вычисление $f$ на
итерацию, коэффициент $1-\alpha \approx 0{,}618$.

\paragraph{Метод Фибоначчи.}
$F_0 = F_1 = 1$, $F_{i+1} = F_i + F_{i-1}$; число шагов $n$ фиксируется
заранее и выбирается из требуемой точности: $F_{n} \ge (b-a)/\delta$. На
$k$-м шаге
\[
  \lambda_k = a_k + \frac{F_{n-k-1}}{F_{n-k+1}}L_k ,
  \qquad
  \mu_k = a_k + \frac{F_{n-k}}{F_{n-k+1}}L_k ,
\]
правило сокращения прежнее, и так же переиспользуется одна точка. Итоговая
длина интервала $(b-a)/F_n$ — \textbf{наилучшая возможная} при заданном числе
вычислений $f$, то есть метод оптимален. Золотое сечение — его предельный
вариант: $F_{n-1}/F_n \to 0{,}618$, но ему не нужно знать $n$ заранее.

\paragraph{Сводка (коэффициент сокращения на одно вычисление $f$).}
равномерный поиск $\approx 0{,}794$; дихотомия и бисекция $\approx 0{,}707$;
золотое сечение $\approx 0{,}618$; Фибоначчи — оптимум.

\seefull{гл.~2, \S9--\S10 (численный пример для всех шести методов там же)}

\section{Билет 6. Градиентные методы первого порядка}

\ticketq{Градиентные методы первого порядка с различными способами выбора
шага.}

\paragraph{Постановка.}
$f(x) \to \min$, $x \in \R^n$, $f$ дифференцируема. Градиент $\nabla f(x)$ —
направление наискорейшего возрастания, $-\nabla f(x)$ — антиградиент. Схема
\[
  x_{k+1} = x_k - \alpha_k \nabla f(x_k) , \qquad \alpha_k > 0 .
\]
Останов: $\norm{\nabla f(x_k)} < \eps$.

\textbf{Направление спуска} $d_k$ в точке $x_k$: $\nabla\T f(x_k)d_k < 0$;
тогда $f(x_k + \alpha d_k) < f(x_k)$ при достаточно малом $\alpha > 0$.

\paragraph{Теорема (антиградиент — направление спуска).}
\textit{$\nabla f(x_k) \ne 0$ $\Rightarrow$ $d_k = -\nabla f(x_k)$ —
направление спуска.}

\textsc{Док-во.}
$\nabla\T f(x_k)d_k = -\nabla\T f(x_k)\nabla f(x_k)
 = -\norm{\nabla f(x_k)}^2 < 0$. $\blacksquare$

\paragraph{Лемма (основная оценка).}
\textit{Пусть $\nabla f$ липшицев: $\norm{\nabla f(x) - \nabla f(y)} \le
R\norm{x-y}$. Тогда $\forall \alpha > 0$}
\[
  f\big(x_k - \alpha\nabla f(x_k)\big) - f(x_k)
  \ \le\ -\alpha(1 - \alpha R)\,\norm{\nabla f(x_k)}^2 .
\]

\textsc{Док-во.} $x = x_k - \alpha\nabla f(x_k)$, $x - x_k =
-\alpha\nabla f(x_k)$. По формуле конечных приращений и прибавив-вычтя
$\nabla\T f(x_k)(x-x_k)$:
\[
  f(x) - f(x_k)
  = \underbrace{\nabla\T f(x_k)(x-x_k)}_{=\,-\alpha\norm{\nabla f(x_k)}^2}
  + \big(\nabla f(x_{\text{ср}}) - \nabla f(x_k)\big)\T(x - x_k) .
\]
Второе слагаемое — по Коши--Буняковскому и Липшицу, с учётом
$\norm{x_{\text{ср}} - x_k} \le \norm{x-x_k}$:
\[
  \le R\norm{x-x_k}^2 = R\alpha^2\norm{\nabla f(x_k)}^2
  \ \Rightarrow\
  f(x) - f(x_k) \le \big(R\alpha^2 - \alpha\big)\norm{\nabla f(x_k)}^2 .
  \ \blacksquare
\]

\paragraph{Вспомогательные неравенства (сильная выпуклость).}
При $m\norm{z}^2 \le z\T H(y)z \le M\norm{z}^2$ из
$f(y) \ge f(x) + g\T(y-x) + \tfrac{m}{2}\norm{y-x}^2$ минимизацией правой
части по $y$ получаем $f_* \ge f(x) - \tfrac{1}{2m}\norm{g}^2$, то есть
\[
  \norm{g}^2 \ \ge\ 2m\big(f(x) - f_*\big) ,
  \qquad
  \tfrac{m}{2}\norm{x-x_*}^2 \le f(x) - f_* \le \tfrac{M}{2}\norm{x-x_*}^2 ,
\]
и $\nabla f$ липшицев с $R = M$. Эти три оценки — рабочий инструмент всех
теорем о скорости: первая переводит $\norm{g}^2$ в убывание по значению
функции, вторая — оценку по значению в оценку по аргументу.

\paragraph{1) Постоянный шаг.}
$\alpha_k \equiv \alpha = \mathrm{const}$.

\textbf{Теорема.} \textit{Пусть $m\norm{z}^2 \le z\T H(y)z \le M\norm{z}^2$,
$0 < m \le M$, и $0 < \alpha < 2/M$. Тогда}
\[
  \norm{x_{k+1} - x_*} \le q\,\norm{x_k - x_*} ,
  \qquad q = \max\{\abs{1-\alpha m},\ \abs{1-\alpha M}\} ,
\]
\[
  q_{\min} = \frac{M-m}{M+m}
  \qquad\text{при}\qquad \alpha = \frac{2}{M+m} .
\]

\textsc{Док-во.} $\nabla f(x_*) = 0$, по формуле конечных приращений
$\nabla f(x_k) - \nabla f(x_*) = H(\bar{x}_k)(x_k - x_*)$, поэтому
\[
  x_{k+1} - x_* = \big(E - \alpha H(\bar{x}_k)\big)(x_k - x_*)
  \ \Rightarrow\
  \norm{x_{k+1} - x_*} \le \norm{E - \alpha H(\bar{x}_k)}\cdot\norm{x_k - x_*} .
\]
$H$ симметрична, её собственные числа в $[m,M]$ $\Rightarrow$ у
$E - \alpha H$ они равны $1 - \alpha\mu$, $\mu \in [m,M]$, и
\[
  \norm{E - \alpha H} = \max_{\mu \in [m,M]}\abs{1-\alpha\mu}
   = \max\{\abs{1-\alpha m}, \abs{1-\alpha M}\} = q
\]
(модуль линейной функции максимален на конце отрезка). Минимум по $\alpha$ —
там, где ветви сравниваются: $1 - \alpha m = \alpha M - 1$ $\Rightarrow$
$\alpha = 2/(m+M)$, $q_{\min} = 1 - \frac{2m}{m+M} = \frac{M-m}{M+m}$.
Отсюда же оценка по значению функции:
\[
  f(x_k) - f(x_*) \ \le\ \tfrac{M}{2}\norm{x_k-x_*}^2 \ \le\ C q^{2k} ,
  \qquad C = \tfrac{M}{2}\norm{x_0 - x_*}^2 . \ \blacksquare
\]

Шаг надо знать заранее; при $\alpha \ge 2/M$ метод расходится.

\paragraph{2) Дробление шага (правило Армихо).}
$\alpha_0 > 0$, $\lambda \in (0,1)$, $\eps \in (0,1)$. Шаг принимается при
\[
  f\big(x_k - \alpha\nabla f(x_k)\big) - f(x_k)
  \ \le\ -\eps\,\alpha\,\norm{\nabla f(x_k)}^2 ,
\]
иначе $\alpha := \lambda\alpha$.

\textbf{Конечность дробления.} По лемме условие заведомо выполнено, если
\[
  -\alpha(1-\alpha R) \le -\eps\alpha
  \ \iff\ 1 - \alpha R \ge \eps
  \ \iff\ \alpha \le \frac{1-\eps}{R} ,
\]
то есть при достаточно малом $\alpha$ $\Rightarrow$ дробление конечно и шаг
ограничен снизу.

\textbf{Теорема (сходимость).} \textit{$f \in C^1$, ограничена снизу,
$\nabla f$ липшицев $\Rightarrow$ $\norm{\nabla f(x_k)} \to 0$ при любом
$x_0$.} Если добавить сильную выпуклость ($m, M$ как выше), то
$x_k \to x_*$ и
$f(x_{k+1}) - f(x_*) \le q\big(f(x_k) - f(x_*)\big)$, $q \in (0,1)$ —
линейная скорость.

\paragraph{3) Наискорейший спуск.}
$\alpha_k = \argmin_{\alpha > 0} f\big(x_k - \alpha\nabla f(x_k)\big)$
(одномерная задача решается методами билета~5).

\textbf{Теорема (сходимость).} При условиях предыдущей теоремы
$\norm{\nabla f(x_k)} \to 0$: наискорейший спуск не хуже любого
фиксированного $\alpha$, и дальше работает та же лемма.

\textbf{Теорема (ортогональность соседних градиентов).}
\textit{Если $\alpha_k$ найден точно и минимум внутренний, то}
$\nabla\T f(x_{k+1})\nabla f(x_k) = 0$.

\textsc{Док-во.} $\varphi(\alpha) = f(x_k - \alpha g_k)$, $g_k = \nabla f(x_k)$;
внутренний минимум $\Rightarrow$ $\varphi'(\alpha_k) = 0$, а
\[
  \varphi'(\alpha) = \nabla\T f(x_k - \alpha g_k)(-g_k)
  \ \Rightarrow\
  0 = -\nabla\T f(x_{k+1})g_k . \ \blacksquare
\]
Отсюда траектория — зигзаг с прямыми углами; при $M \gg m$ (вытянутые линии
уровня) $q \to 1$ и сходимость медленная.

\paragraph{Сравнение.}
Постоянный шаг — дёшево, но $\alpha$ надо подобрать заранее по $m, M$;
дробление — не требует знания констант, шаг подстраивается; наискорейший
спуск — лучший шаг вдоль направления ценой одномерной минимизации на каждой
итерации. У всех трёх скорость линейная, с $q$, зависящим от
обусловленности $m/M$.

\seefull{гл.~2, \S12--\S16}

\section{Билет 7. Градиентные методы второго порядка}

\ticketq{Градиентные методы второго порядка с различными способами выбора
шага.}

\paragraph{Постановка и вывод направления.}
$f(x) \to \min$, $x \in \R^n$, $f \in C^2$; $g_k = \nabla f(x_k)$,
$H_k = \nabla^2 f(x_k)$. В окрестности $x_k$ функция заменяется квадратичной
моделью ($p = x - x_k$):
\[
  \psi_k(x_k + p) = f(x_k) + g_k\T p + \tfrac12 p\T H_k p .
\]
Необходимое условие минимума модели:
\[
  \nabla_p \psi_k = g_k + H_k p = 0
  \ \Rightarrow\ H_k p_k = -g_k
  \ \Rightarrow\ p_k = -H_k^{-1} g_k ,
  \qquad x_{k+1} = x_k + \alpha_k p_k .
\]

\paragraph{Теорема (направление Ньютона — направление спуска).}
\textit{$H_k \succ 0$, $g_k \ne 0$ $\Rightarrow$ $g_k\T p_k < 0$.}

\textsc{Док-во.} $g_k\T p_k = -g_k\T H_k^{-1}g_k$; из $H_k \succ 0$ следует
$H_k^{-1} \succ 0$ $\Rightarrow$ $g_k\T H_k^{-1}g_k > 0$ $\Rightarrow$
$g_k\T p_k < 0$. В количественной форме, подставив $g_k = -H_kp_k$:
\[
  g_k\T p_k = -p_k\T H_k p_k \ \le\ -m\norm{p_k}^2 < 0 . \ \blacksquare
\]

\paragraph{Теорема (квадратичная задача за одну итерацию).}
\textit{Если $f(x) = \tfrac12 x\T A x + b\T x + c$, $A = A\T \succ 0$, то
метод Ньютона с $\alpha_k = 1$ находит минимум за один шаг.}

\textsc{Док-во.} $\nabla f = Ax + b$, $H = A$ $\Rightarrow$
$p_k = -A^{-1}(Ax_k + b)$, откуда
\[
  x_{k+1} = x_k - A^{-1}Ax_k - A^{-1}b = -A^{-1}b ,
\]
а из $\nabla f(x_*) = Ax_* + b = 0$ следует $x_* = -A^{-1}b = x_{k+1}$.
$\blacksquare$

\paragraph{Способы выбора шага.}
\begin{csteps}
  \item $\alpha_k \equiv 1$ — \textbf{чистый метод Ньютона}. Быстр вблизи
        $x_*$, но вдали может увеличить $f$.
  \item \textbf{Дробление}: $\alpha_0 = 1$, $\lambda \in (0,1)$,
        $\eps \in (0, \tfrac12]$, условие приёма
        \[
          f(x_k + \alpha p_k) - f(x_k) \le \alpha\,\eps\,\nabla\T f(x_k)p_k ,
        \]
        иначе $\alpha := \lambda\alpha$ — \textbf{демпфированный Ньютон}.
  \item \textbf{Наискорейший спуск вдоль $p_k$}:
        $\alpha_k = \argmin_{\alpha>0} f(x_k + \alpha p_k)$.
\end{csteps}

\textbf{Почему $\eps \le \tfrac12$.} Для квадратичной $f$ вдоль
ньютоновского направления $f(x_k + \alpha p) - f(x_k) =
\nabla\T f\,p\,(\alpha - \alpha^2/2)$, при $\alpha = 1$ это
$\tfrac12\nabla\T f\,p$. Условие приёма даёт
\[
  \tfrac12\nabla\T f\,p \le \eps\,\nabla\T f\,p
  \ \iff\ \eps \le \tfrac12
\]
(неравенство переворачивается, так как $\nabla\T f\,p < 0$). При
$\eps > \tfrac12$ дробление отвергало бы идеальный единичный шаг и метод терял
бы квадратичную скорость.

\paragraph{Теорема (квадратичная сходимость).}
\textit{Пусть}
\textit{1) $m\norm{z}^2 \le z\T H(y)z \le M\norm{z}^2$, $0 < m \le M$;}
\textit{2) $\norm{H(x) - H(y)} \le R\norm{x-y}$,}
\textit{и шаг выбирается дроблением. Тогда $x_k \to x_*$ при любом $x_0$,
причём}
\[
  \norm{x_{k+1} - x_*} \ \le\ \frac{R}{m}\,\norm{x_k - x_*}^2 .
\]
\textsc{Док-во (глобальная часть).} Из $\norm{H} \le M$:
$f(x_k+\alpha p_k) - f(x_k) \le \alpha g_k\T p_k + \tfrac{M}{2}\alpha^2\norm{p_k}^2$,
а $g_k\T p_k \le -m\norm{p_k}^2$. Условие дробления выполнено, как только
$\tfrac{M}{2}\alpha^2\norm{p_k}^2 \le (1-\eps)\alpha\,\abs{g_k\T p_k}$, то есть
при $\alpha \le 2(1-\eps)m/M$ $\Rightarrow$ дробление конечно и
$\alpha_k \ge \bar\alpha > 0$. Тогда, учитывая
$\norm{p_k} = \norm{H_k^{-1}g_k} \ge \norm{g_k}/M$,
\[
  f(x_k) - f(x_{k+1}) \ \ge\ \eps\bar\alpha m \norm{p_k}^2
   \ \ge\ \frac{\eps\bar\alpha m}{M^2}\,\norm{g_k}^2 ,
\]
и дальше как в билете~6: $f(x_k)$ убывает и ограничена снизу $\Rightarrow$
$\sum_k\norm{g_k}^2 < \infty$ $\Rightarrow$ $\norm{g_k}\to0$, а по сильной
выпуклости $x_k \to x_*$.

\textsc{Док-во (квадратичная часть).} $e_k = x_k - x_*$, $\alpha_k = 1$,
$\nabla f(x_*) = 0$. Пользуясь $g_k = \int_0^1 H(x_*+te_k)e_k\,dt$:
\[
  e_{k+1} = e_k - H_k^{-1}g_k
   = H_k^{-1}\int_0^1 \big(H_k - H(x_*+te_k)\big)e_k\,dt ,
\]
а из липшицевости $\norm{H_k - H(x_*+te_k)} \le R(1-t)\norm{e_k}$ и
$\int_0^1(1-t)dt = \tfrac12$ получаем
$\norm{e_{k+1}} \le \tfrac{R}{2m}\norm{e_k}^2$. $\blacksquare$

Условие~1) даёт существование, единственность $x_*$ и обратимость $H$
(направление определено и является спуском); условие~2) — липшицевость
\emph{гессиана} — отвечает за квадратичную оценку. Глобальность обеспечивает
дробление: вдали от $x_*$ шаг уменьшается и $f$ убывает, вблизи единичный шаг
проходит проверку при $\eps \le \tfrac12$.

\paragraph{Модификации (экономия на гессиане).}
\begin{csteps}
  \item \textbf{Фиксированная матрица:} $H_0 = \nabla^2 f(x_0)$,
        $x_{k+1} = x_k - \alpha_k H_0^{-1}\nabla f(x_k)$ — одно разложение на
        весь счёт.
  \item \textbf{Периодическое обновление:} гессиан фиксируется на $t$ итераций,
        \[
          x_{lt+i+1} = x_{lt+i} - \alpha_{lt+i}H^{-1}(x_{lt})\nabla f(x_{lt+i}),
          \quad i = \overline{0,t-1} .
        \]
        $t = 1$ — обычный Ньютон, $t = \infty$ — фиксированная матрица.
  \item \textbf{Левенберг--Марквардт:} $(H_k + \lambda_k E)p_k = -g_k$.
        $\lambda_k \to 0$ — Ньютон; $\lambda_k \to \infty$ —
        $p_k \approx -\frac{1}{\lambda_k}g_k$, градиентный шаг малой длины.
        Слагаемое $\lambda_k E$ делает матрицу положительно определённой даже
        там, где $H_k$ таковой не является.
\end{csteps}
Модификации 1--2 сохраняют сходимость (направление остаётся спуском, пока
матрица $\succ 0$), но теряют квадратичную скорость: оценка опиралась на то,
что модель построена \emph{в текущей} точке.

\paragraph{Ограничения методов второго порядка.}
Нужно вычислять $H_k$ и решать $H_k p_k = -g_k$ — дорого при большом $n$;
при $H_k \not\succ 0$ направление может не быть спуском, и дробление этого не
исправляет; при плохом $x_0$ чистый метод может расходиться.

\seefull{гл.~2, \S20--\S22 и \S24}

\section{Билет 8. Метод Хука--Дживса}

\ticketq{Метод Хука--Дживса.}

\paragraph{Постановка.}
$f(x) \to \min$, $x \in \R^n$. Метод \textbf{прямой}, или \textbf{нулевого
порядка}: использует только значения $f(x_1), f(x_2), \dots$ и не требует
$\nabla f(x)$, $H(x) = \nabla^2 f(x)$. Строится последовательность
$x^0, x^1, \dots$ с $f(x^{k+1}) < f(x^k)$ по возможности.

\textbf{Параметры:} начальная точка $x^0$; шаги по координатам $h_i > 0$,
$i = \overline{1,n}$; точность $\eps > 0$; коэффициент дробления
$\beta \in (0,1)$. \textbf{Базисная точка} $x^b$ — текущая основная точка,
вначале $x^b = x^0$.

Метод состоит из двух чередующихся шагов: исследующего поиска и поиска по
образцу.

\paragraph{Шаг 1. Исследующий поиск.}
Покоординатная проверка из текущей точки $x$; $e_i$ — $i$-й единичный вектор.
Для каждой координаты $i = \overline{1,n}$:
\[
  \begin{aligned}
    f(x + h_i e_i) &< f(x) &&\Rightarrow\ x := x + h_i e_i ; \\
    \text{иначе } f(x - h_i e_i) &< f(x) &&\Rightarrow\ x := x - h_i e_i ; \\
    \text{иначе} &&&\ \text{координата не меняется.}
  \end{aligned}
\]
Сначала пробуем «вперёд», потом «назад». После всех координат получается
точка $x^e$ — результат исследующего поиска. За один проход каждая координата
меняется не более одного раза, и $f$ только невозрастает.

\paragraph{Шаг 2. Поиск по образцу.}
Если $f(x^e) < f(x^b)$, поиск удачен; направление удачного смещения
$x^e - x^b$, и строится ускоренная точка
\[
  x^p = x^e + (x^e - x^b) = 2x^e - x^b .
\]
Из $x^p$ выполняется \emph{ещё один исследующий поиск}, дающий $x^{e2}$, и
сравнение идёт уже по нему:
\[
  f(x^{e2}) < f(x^e) \ \Rightarrow\ x^b := x^e,\ x^e := x^{e2},\
    \text{повторить поиск по образцу;}
\]
\[
  f(x^{e2}) \ge f(x^e) \ \Rightarrow\ x^b := x^e,\
    \text{вернуться к исследующему поиску.}
\]
Смысл: исследующий поиск за цикл выявил удачное направление — есть смысл
пройти по нему ещё раз, не тратя вычислений на пробные шаги. Это дешёвая
замена одномерной минимизации. Повторный исследующий поиск из $x^p$ нужен
затем, чтобы подправить направление: чистая экстраполяция легко проскакивает
мимо дна оврага.

\paragraph{Условие остановки.}
Если исследующий поиск не дал улучшения, шаги дробятся: $h_i := \beta h_i$.
Метод останавливается при $\max_i h_i < \eps$ — поиск с достаточно мелкой
сеткой уже не улучшает значение.

\paragraph{Теорема (о монотонном убывании).}
\textit{Если новая базисная точка принимается только при строгом уменьшении
$f$, то $f(x^b_{k+1}) \le f(x^b_k)$, а при удачной итерации
$f(x^b_{k+1}) < f(x^b_k)$.}

\textsc{Док-во.} Два случая.

\emph{1) Поиск неудачен:} $f(x^e) \ge f(x^b_k)$ $\Rightarrow$ базисная точка
не меняется, $x^b_{k+1} = x^b_k$ $\Rightarrow$ $f(x^b_{k+1}) = f(x^b_k)$
(дробятся лишь шаги $h_i$).

\emph{2) Поиск удачен:} $f(x^e) < f(x^b_k)$, строится $x^p = 2x^e - x^b_k$.
\[
  f(x^p) < f(x^e) \Rightarrow x^b_{k+1} = x^p :\ f(x^p) < f(x^e) < f(x^b_k) ;
\]
\[
  f(x^p) \ge f(x^e) \Rightarrow x^b_{k+1} = x^e :\ f(x^e) < f(x^b_k) .
\]
В обоих подслучаях $f(x^b_{k+1}) < f(x^b_k)$. $\blacksquare$

\paragraph{Теорема (о конечности исследующего поиска).}
\textit{Один исследующий поиск требует не более $2n$ вычислений $f$.}

\textsc{Док-во.} Для каждой координаты проверяются не более двух точек,
$x + h_i e_i$ и $x - h_i e_i$ $\Rightarrow$ не более двух вычислений на
координату; координат $n$ $\Rightarrow$ не более $2n$. Так как $n$ конечно,
поиск завершается за конечное число операций. $\blacksquare$

\paragraph{Алгоритм.}
\begin{csteps}
  \item Задать $x^0$, $h_i > 0$, $\eps > 0$, $\beta \in (0,1)$; $x^b := x^0$.
  \item Выполнить исследующий поиск из $x^b$, получить $x^e$.
  \item $f(x^e) \ge f(x^b)$: положить $h_i := \beta h_i$; если
        $\max_i h_i < \eps$ — конец, иначе к~шагу~2.
  \item $f(x^e) < f(x^b)$: построить $x^p = 2x^e - x^b$ и выполнить
        исследующий поиск из $x^p$, получив $x^{e2}$.
  \item $f(x^{e2}) < f(x^e)$: положить $x^b := x^e$, $x^e := x^{e2}$ и
        повторить шаг~4. Иначе $x^b := x^e$ и к~шагу~2.
\end{csteps}

\paragraph{Сходимость и ограничения.}
Метод не требует дифференцируемости и применим, когда производные неизвестны
или дороги. Глобального минимума для невыпуклой $f$ он не гарантирует —
сходится к локальному минимуму либо к точке, где выбранные координатные шаги
уже не дают улучшения. Прочие ограничения: много вычислений $f$; скорость ниже
градиентных и ньютоновских методов; сильная зависимость от выбора $h_i$ и от
масштаба переменных (при разных масштабах метод работает медленно); крупные
шаги проскакивают хорошую область, мелкие замедляют сходимость.

\seefull{гл.~2, \S19 (численный пример там же)}

\section{Билет 9. Квазиньютоновские методы}

\ticketq{Квазиньютоновские методы.}

\paragraph{Идея.}
Метод Ньютона требует на каждой итерации вычислять $H(x_k) = \nabla^2 f(x_k)$ и
решать $H_k p_k = -g_k$ — дорого. Квазиньютоновские методы гессиан не
вычисляют: вместо $H^{-1}(x_k)$ берётся приближение $A_k$,
\[
  p_k = -A_k \nabla f(x_k) ,
  \qquad
  x_{k+1} = x_k + \alpha_k p_k ,
  \qquad
  A_k \xrightarrow[k\to\infty]{} H^{-1}(x_*) ,
\]
причём $A_{k+1} = A_k + \Delta A_k$ и $A_1 = E$ (первая итерация —
градиентный шаг, информации о кривизне ещё нет). Методы занимают промежуточное
положение: используют информацию первого порядка, но имитируют поведение
второго.

\paragraph{Разности и секущее условие.}
\[
  s_k = x_{k+1} - x_k , \qquad
  y_k = \nabla f(x_{k+1}) - \nabla f(x_k) .
\]
Для точного гессиана $y_k \approx H_{k+1}s_k$, поэтому от приближения обратной
матрицы требуют \textbf{секущее (квазиньютоновское) условие}
\[
  A_{k+1}\,y_k = s_k ,
\]
то есть чтобы $A_{k+1}$ действовала на уже известном направлении так же, как
настоящая $H^{-1}$. \textbf{Условие кривизны:} $s_k\T y_k > 0$.

\paragraph{Теорема (направление спуска).}
\textit{$A_k \succ 0$, $g_k \ne 0$ $\Rightarrow$ $p_k = -A_k g_k$ —
направление спуска.}

\textsc{Док-во.} $g_k\T p_k = -g_k\T A_k g_k < 0$, так как $A_k \succ 0$ и
$g_k \ne 0$. $\blacksquare$

Поэтому положительная определённость $A_k$ — то, что требуется сохранять при
обновлении.

\paragraph{Формула ДФП (Дэвидон--Флетчер--Пауэлл).}
\[
  A_{k+1} = A_k + \frac{s_k s_k\T}{s_k\T y_k}
            - \frac{A_k y_k y_k\T A_k}{y_k\T A_k y_k} .
\]

\textbf{Теорема (секущее условие).} \textit{$A_{k+1}y_k = s_k$.}

\textsc{Док-во.} Умножим формулу справа на $y_k$:
\[
  A_{k+1}y_k = A_k y_k
    + \frac{s_k\,(s_k\T y_k)}{s_k\T y_k}
    - \frac{A_k y_k\,(y_k\T A_k y_k)}{y_k\T A_k y_k}
    = A_k y_k + s_k - A_k y_k = s_k ,
\]
так как $s_k\T y_k$ и $y_k\T A_k y_k$ — \emph{числа}. $\blacksquare$

\textbf{Теорема (положительная определённость).}
\textit{$A_k \succ 0$, $s_k\T y_k > 0$ $\Rightarrow$ $A_{k+1} \succ 0$.}

\textsc{Док-во.} Для $z \ne 0$
\[
  z\T A_{k+1} z = \underbrace{z\T A_k z
    - \frac{(z\T A_k y_k)^2}{y_k\T A_k y_k}}_{\ \ge\ 0}
    + \underbrace{\frac{(z\T s_k)^2}{s_k\T y_k}}_{\ \ge\ 0} \ \ge\ 0 ,
\]
первое слагаемое неотрицательно по Коши--Буняковскому в метрике
$\langle u,v\rangle_{A_k} = u\T A_k v$, второе — из $s_k\T y_k > 0$.
Равенство нулю требует обнуления обоих: из равенства в Коши следует
$z = c\,y_k$, а из $(z\T s_k)^2 = 0$ тогда $c\,y_k\T s_k = 0$. Но
$y_k\T s_k > 0$, значит $c = 0$ и $z = 0$. Поэтому при $z \ne 0$ строго
$z\T A_{k+1}z > 0$. $\blacksquare$

\paragraph{Формула БФГШ (Бройден--Флетчер--Голдфарб--Шенно).}
При $\rho_k = \dfrac{1}{y_k\T s_k}$
\[
  A_{k+1} = \big(E - \rho_k s_k y_k\T\big) A_k
            \big(E - \rho_k y_k s_k\T\big) + \rho_k s_k s_k\T .
\]

\textbf{Теорема (секущее условие).} \textit{$A_{k+1}y_k = s_k$.}

\textsc{Док-во.} $\rho_k(s_k\T y_k) = 1$ $\Rightarrow$
$(E - \rho_k y_k s_k\T)y_k = y_k - \rho_k y_k (s_k\T y_k) = 0$ $\Rightarrow$
первое слагаемое обнуляется. Остаётся
$\rho_k s_k s_k\T y_k = \frac{1}{y_k\T s_k}s_k (s_k\T y_k) = s_k$.
$\blacksquare$

\textbf{Теорема (положительная определённость).}
\textit{$A_k \succ 0$, $y_k\T s_k > 0$ $\Rightarrow$ $A_{k+1} \succ 0$.}

\textsc{Док-во.} Обозначим $w = (E - \rho_k y_k s_k\T)z$, тогда
\[
  z\T A_{k+1} z = w\T A_k w + \rho_k (z\T s_k)^2 \ \ge\ 0 .
\]
Равенство $\Rightarrow$ $w\T A_k w = 0$ и $(z\T s_k)^2 = 0$. Из
$A_k \succ 0$ следует $w = 0$; из $s_k\T z = 0$ получаем
$w = z - \rho_k y_k(s_k\T z) = z$ $\Rightarrow$ $z = 0$. $\blacksquare$

\paragraph{Алгоритм.}
\begin{csteps}
  \item Задать $x_0$, $\eps > 0$, $A_0 = E$, $k = 0$.
  \item $g_k = \nabla f(x_k)$; если $\norm{g_k} < \eps$ — конец,
        $x_* \approx x_k$.
  \item $p_k = -A_k g_k$.
  \item $\alpha_k = \argmin_{\alpha>0} f(x_k + \alpha p_k)$ (либо условия
        Армихо/Вульфа); $x_{k+1} = x_k + \alpha_k p_k$.
  \item $s_k = x_{k+1} - x_k$, $y_k = \nabla f(x_{k+1}) - \nabla f(x_k)$.
  \item Если $k$ принадлежит заданному множеству номеров обновления $\Pi_0$ —
        \textbf{рестарт}: $A_{k+1} = E$ (накопленная матрица сбрасывается,
        обычно раз в $n$ итераций). Иначе при $s_k\T y_k > 0$ обновить $A_k$
        по ДФП или БФГШ; при $s_k\T y_k \le 0$ — тоже сброс $A_{k+1} = E$.
        $k := k+1$, к~шагу~2.
\end{csteps}

\paragraph{Свойства и ограничения.}
Условие кривизны $s_k\T y_k > 0$ выполняется автоматически для сильно выпуклой
$f$ при точном одномерном поиске; в общем случае его обеспечивают условиями
Вульфа. БФГШ устойчивее ДФП: обновление записано как сопряжение
$W A_k W\T$ плюс слагаемое ранга один, и положительная определённость не
теряется от накопления ошибок округления, тогда как в ДФП она держится на
разности двух слагаемых. Для квадратичной функции $n$ переменных при точном
одномерном поиске метод сходится не более чем за $n$ шагов, и $A_n = H^{-1}$.

Ограничения: нужно хранить и обновлять матрицу $n \times n$ (при большом $n$
используют L-BFGS); чувствительность к точности одномерного поиска; для
невыпуклых функций — только локальный минимум.

\seefull{гл.~2, \S23 (численный пример ДФП там же)}

\section{Билет 10. Метод возможных направлений}

\ticketq{Метод возможных направлений.}

\paragraph{Постановка.}
\[
  \varphi_0(x) \to \min_{x \in S} ,
  \qquad
  S = \{\, x \in \R^n \mid \varphi_i(x) \le 0,\ i = \overline{1,m};\ Ax = b \,\} ,
\]
Метод строит допустимые точки $x_k \in S$ с убыванием цели по схеме
$x_{k+1} = x_k + \alpha_k s_k$, где $s_k$ — возможное направление спуска,
$\alpha_k > 0$ — шаг. Требуется допустимая начальная точка $x_0 \in S$.

\textbf{Предположения:} $\varphi_i$ выпуклы и дифференцируемы, градиенты
ограничены и липшицевы,
\[
  \norm{\nabla\varphi_i(x)} \le K , \qquad
  \norm{\nabla\varphi_i(x) - \nabla\varphi_i(y)} \le R\norm{x-y} ,
\]
$S$ компактно, выполнено условие невырожденности.

\paragraph{Определения.}
Ненулевой $s$ — \textbf{возможное направление} в $x \in S$, если
$\exists\,\delta > 0 : x + \alpha s \in S$ при всех $0 < \alpha < \delta$.
Если дополнительно $\varphi_0(x + \alpha s) < \varphi_0(x)$, то $s$ —
\textbf{возможное направление спуска}.

Ограничения в точке $x$: \textbf{активные}
$I_0(x) = \{i \mid \varphi_i(x) = 0\}$; \textbf{пассивные}
$\{i \mid \varphi_i(x) < 0\}$; \textbf{почти активные}
$I_\delta(x) = \{i \mid -\delta \le \varphi_i(x) \le 0\}$ — учитываются,
чтобы точка не подходила вплотную к границе.

\paragraph{Лемма (о возможном направлении).}
\textit{Пусть $x \in S$ и}
\[
  \nabla\T\varphi_i(x)\,s < 0 \ \ \forall i \in I_0(x) ,
  \qquad As = 0 .
\]
\textit{Тогда $s$ — возможное направление в $x$.}

\textsc{Док-во.} \emph{Равенства:} $Ax = b$, $As = 0$ $\Rightarrow$
$A(x + \alpha s) = b + \alpha\cdot 0 = b$ при любом $\alpha$.

\emph{Активные $i \in I_0(x)$:} $\varphi_i(x) = 0$, по формуле Тейлора
\[
  \varphi_i(x + \alpha s) = \alpha\,\nabla\T\varphi_i(x)s + o(\alpha) ,
\]
главный член отрицателен $\Rightarrow$ $\varphi_i(x+\alpha s) < 0$ при малом
$\alpha > 0$.

\emph{Пассивные:} $\varphi_i(x) < 0$ и $\varphi_i$ непрерывна $\Rightarrow$
$\varphi_i(x + \alpha s) < 0$ при малом $\alpha$.

Значит $x + \alpha s \in S$. $\blacksquare$

\paragraph{Лемма (о направлении спуска).}
\textit{$\nabla\T\varphi_0(x)s < 0$ $\Rightarrow$ $s$ — направление спуска.}

\textsc{Док-во.}
$\varphi_0(x+\alpha s) - \varphi_0(x) = \alpha\nabla\T\varphi_0(x)s + o(\alpha)$;
главный линейный член отрицателен $\Rightarrow$ при малом $\alpha > 0$ разность
отрицательна. $\blacksquare$

Итого: одновременное выполнение $\nabla\T\varphi_0(x)s < 0$,
$\nabla\T\varphi_i(x)s < 0$ ($i \in I_0$), $As = 0$ даёт возможное направление
спуска.

\paragraph{Вспомогательная задача Зойтендейка.}
Направление ищется как решение задачи ЛП: при $\xi_i > 0$, $i = \overline{0,m}$,
\[
  \eta \to \min_{\mathbb{H}} ,
  \quad
  \mathbb{H} = \Big\{ \eta \ \Big|\
    \nabla\T\varphi_i(x)s \le \xi_i\eta,\ i \in I_\delta(x)\cup\{0\} ;\
    As = 0 ;\ \norm{s}_\infty \le 1 \Big\} .
\]
Пара $s = 0$, $\eta = 0$ допустима $\Rightarrow$ $\eta_\delta \le 0$.

\textbf{Смысл.} При $\eta < 0$ и $\xi_i > 0$ имеем $\xi_i\eta < 0$, то есть
$\nabla\T\varphi_0 s < 0$ и $\nabla\T\varphi_i s < 0$ — по леммам $s$ есть
возможное направление спуска. Условие $As = 0$ сохраняет равенства, а
$\norm{s}_\infty \le 1$ нормирует задачу: без него $s$ можно было бы
растягивать и $\eta \to -\infty$. Случай $\eta_\delta = 0$ означает, что
возможного направления спуска нет, то есть точка приближённо стационарна.

\paragraph{Лемма (оценки вдоль направления).}
\textit{При $\alpha > 0$, $\norm{s_\delta} \le 1$:}
\[
  \varphi_0(x + \alpha s_\delta) \le \varphi_0(x) + \alpha\xi_0\eta_\delta
    + \alpha^2 R\norm{s_\delta}^2 ,
  \quad
  \varphi_i(x + \alpha s_\delta) \le \varphi_i(x)
    + \alpha\nabla\T\varphi_i(x)s_\delta + \alpha^2 R\norm{s_\delta}^2 ,
\]
\textit{и $A(x+\alpha s_\delta) = b$.} Первое — из липшицевости градиента и
ограничения $\nabla\T\varphi_0 s \le \xi_0\eta_\delta$.

\paragraph{Допустимость и убывание при $\eta_\delta < 0$.}
Разбор по трём группам даёт явные границы шага:
\begin{csteps}
  \item $i \in I_\delta(x)$ (почти активные): $\varphi_i(x) \le 0$ и
        $\nabla\T\varphi_i s_\delta \le \xi_i\eta_\delta < 0$, поэтому
        $\varphi_i(x+\alpha s_\delta) \le \alpha\xi_i\eta_\delta + \alpha^2 R
         \le 0$ при $\alpha \le \xi_i\abs{\eta_\delta}/R$;
  \item $i \notin I_\delta(x)$ (с запасом): $\varphi_i(x) < -\delta$ и
        $\norm{\nabla\varphi_i} \le K$, поэтому
        $\varphi_i(x+\alpha s_\delta) \le -\delta + \alpha K + \alpha^2 R \le 0$
        при малом $\alpha$;
  \item цель: $\varphi_0(x+\alpha s_\delta) - \varphi_0(x)
        \le \alpha\xi_0\eta_\delta + \alpha^2 R
        \le \tfrac12\alpha\xi_0\eta_\delta$ при
        $\alpha \le \xi_0\abs{\eta_\delta}/(2R)$.
\end{csteps}
Значит дробление $\alpha_k = \alpha_0\lambda^{\,j}$ конечно, шаг ограничен
снизу величиной, зависящей от $\abs{\eta_\delta}$, а $\varphi_0(x_k)$ строго
убывает. При $\eta_\delta = 0$ допустимого направления спуска нет —
выполняются условия оптимальности, и $\delta$ уменьшают. $\blacksquare$

\paragraph{Алгоритм.}
Задать $x_0 \in S$, $\eps > 0$, $\delta_0 > 0$, $\lambda \in (0,1)$.
\begin{csteps}
  \item Найти $I_0(x_k)$ и $I_{\delta_k}(x_k)$.
  \item Решить вспомогательную задачу; получить $s_k$, $\eta_k$.
  \item $\eta_k \ge 0$: направления нет — уменьшить $\delta_k := \lambda\delta_k$
        (при $\delta_k < \eps$ — конец, $x_k$ приближённо стационарна),
        к~шагу~1.
  \item $\eta_k < 0$: найти шаг дроблением — $\alpha_k = \alpha_0\lambda^{\,j}$
        с наименьшим $j$, при котором
        \[
          x_k + \alpha_k s_k \in S
          \quad\text{и}\quad
          \varphi_0(x_k + \alpha_k s_k) - \varphi_0(x_k)
            \le \tfrac12\,\alpha_k\,\xi_0\,\eta_k .
        \]
  \item $x_{k+1} = x_k + \alpha_k s_k$; останов при
        $\norm{x_{k+1}-x_k} < \eps$, иначе $k := k+1$, к~шагу~1.
\end{csteps}

\paragraph{Начальное приближение.}
Если $x_0$ неизвестна, решается вспомогательная задача
\[
  \eta \to \min , \qquad \varphi_i(x) - \eta \le 0,\ i = \overline{1,m},
  \qquad Ax = b .
\]
По условию Слейтера существует $\bar{x}$ с $\varphi_i(\bar{x}) < 0$,
$A\bar{x} = b$ $\Rightarrow$ $\min\eta < 0$, и найденная точка строго внутри
всех неравенств — ровно то, что нужно для старта. Это та же идея, что у
вспомогательной задачи симплекс-метода.

\paragraph{Ограничения.}
Нужна допустимая $x_0$; на каждой итерации решается задача ЛП; при сложных
нелинейных ограничениях трудно выбрать допустимый шаг; без учёта почти
активных ограничений метод неустойчив около границы; для невыпуклых задач —
только локальный минимум.

\seefull{гл.~3, \S1 и \S4 (численный пример в \S4)}

\section{Билет 11. Матрица проектирования. Метод Розена}

\ticketq{Матрица проектирования. Метод проекции градиента Розена.}

\paragraph{Постановка и идея.}
\[
  f(x) \to \min , \qquad Ax = b , \qquad Cx \le d ,
\]
$f$ дифференцируема. В безусловной задаче направление убывания —
$-\nabla f(x_k)$, но оно может вывести за пределы допустимого множества.
Поэтому антиградиент \emph{проектируют} на множество допустимых направлений:
из $A(x + \alpha s) = b$ и $Ax = b$ следует $As = 0$, то есть допустимые
направления лежат в подпространстве $\{s \mid As = 0\}$.

\paragraph{Матрица проектирования.}
$P \in \R^{n\times n}$ — \textbf{матрица проектирования}, если
$P\T = P$ и $P^2 = P$.

\textbf{Лемма (общие свойства).} \textit{$Q = E - P$ — тоже проектор;
$P, Q \succeq 0$; $\operatorname{Im}P \perp \operatorname{Im}Q$; разложение
$x = Px + Qx$ единственно.}

\textsc{Док-во.} $Q\T = E - P\T = Q$ и
$Q^2 = E - 2P + P^2 = E - P = Q$. Далее
$x\T Px = x\T P\T Px = \norm{Px}^2 \ge 0$. Ортогональность:
$(Px_1)\T Qx_2 = x_1\T(P - P^2)x_2 = 0$. Единственность: если
$x = y+z = y_1+z_1$ с $y,y_1 \in L = \operatorname{Im}P$,
$z,z_1 \in L' = \operatorname{Im}Q$, то
$y - y_1 = z_1 - z \in L \cap L' = \{0\}$. $\blacksquare$

Для матрицы активных ограничений $A \in \R^{m\times n}$ с $\rank A = m$
(тогда $AA\T$ невырождена) проектор на $\{s \mid As = 0\}$:
\[
  P = E - A\T\big(AA\T\big)^{-1}A .
\]

\paragraph{Теорема (свойства проектора).}
\textit{1) $P\T = P$; 2) $P^2 = P$; 3) $AP = 0$; 4) $Az = 0 \Rightarrow Pz = z$.}

\textsc{Док-во.}
1) $P\T = E\T - A\T\big((AA\T)^{-1}\big)\T A$; $AA\T$ симметрична
$\Rightarrow$ её обратная симметрична $\Rightarrow$ $P\T = P$.

2) $P^2 = E - 2A\T(AA\T)^{-1}A + A\T(AA\T)^{-1}\underbrace{AA\T(AA\T)^{-1}}_{=\,E}A
   = E - A\T(AA\T)^{-1}A = P$.

3) $AP = A - \underbrace{AA\T(AA\T)^{-1}}_{=\,E}A = A - A = 0$.

4) $Pz = z - A\T(AA\T)^{-1}\underbrace{Az}_{=\,0} = z$. $\blacksquare$

\paragraph{Направление метода.}
$g = \nabla f(x)$, $s = -Pg$.

\textbf{Теорема (направление спуска).} \textit{$Pg \ne 0 \Rightarrow g\T s < 0$.}

\textsc{Док-во.} $P\T = P$, $P^2 = P$ $\Rightarrow$
\[
  g\T s = -g\T P g = -g\T P\T P g = -(Pg)\T(Pg) = -\norm{Pg}^2 < 0 .
  \ \blacksquare
\]

\textbf{Теорема (возможное направление).} \textit{$Ax = b$, $s = -P\nabla f(x)$
$\Rightarrow$ $A(x + \alpha s) = b$.}

\textsc{Док-во.} $As = -AP\nabla f(x) = 0$ по свойству 3 $\Rightarrow$
$A(x+\alpha s) = Ax + \alpha\cdot 0 = b$. $\blacksquare$

\paragraph{Условие оптимальности и множители Лагранжа.}
Пусть $Pg = 0$. Тогда
\[
  g - A\T(AA\T)^{-1}Ag = 0
  \ \Rightarrow\ g = A\T\mu , \quad \mu = (AA\T)^{-1}Ag ,
\]
то есть градиент цели выражается через строки активных ограничений — это
условие Лагранжа $\nabla f(x) + A\T\lambda = 0$ при $\lambda = -\mu$. Для
ограничений вида $a_i\T x \le b_i$ условия Куна--Таккера требуют
$\lambda_i \ge 0$:
\begin{itemize}[nosep,left=0pt]
  \item все $\lambda_i \ge 0$ — точка удовлетворяет необходимым условиям
        оптимальности, конец;
  \item $\exists\,\lambda_i < 0$ — это ограничение исключают из активного
        набора и строят проектор заново.
\end{itemize}

\textbf{Почему исключение работает.} Пусть $u_j < 0$ (множитель при активном
неравенстве $f_j\T x \le g_j$). Уберём строку $f_j$ из $A$, получив $A_1$ и
$P_1$. Тогда $\nabla f = -A_1\T\omega_1 - u_j f_j\T$ и
\[
  P_1\nabla f = -u_j P_1 f_j\T \ne 0
\]
(иначе $u_j = 0$ из линейной независимости строк). Новое направление
$s' = -P_1\nabla f = u_j P_1 f_j\T$ одновременно:
\[
  \nabla\T f\,s' = -\norm{P_1\nabla f}^2 < 0 , \qquad
  A_1 s' = 0 , \qquad
  f_j s' = u_j\norm{P_1 f_j\T}^2 < 0 ,
\]
то есть это направление спуска, сохраняющее остальные активные ограничения и
уводящее \emph{внутрь} по снятому неравенству — оно перестаёт быть активным, и
движение допустимо.

\paragraph{Выбор шага.}
Активный набор в $x_k$: $I_k = \{i \mid a_i\T x_k = b_i\}$, $A_k$ — матрица
его строк, $P_k = E - A_k\T(A_kA_k\T)^{-1}A_k$, $s_k = -P_k g_k$. Для
неактивных ограничений нужно $a_i\T(x_k + \alpha s_k) \le b_i$, откуда
\[
  \alpha_{\max} = \min_{i \notin I_k,\ a_i\T s_k > 0}
    \frac{b_i - a_i\T x_k}{a_i\T s_k}
  \qquad (\alpha_{\max} = +\infty, \text{ если таких } i \text{ нет}) ,
\]
и $\alpha_k = \argmin_{0 \le \alpha \le \alpha_{\max}} f(x_k + \alpha s_k)$.

\paragraph{Алгоритм.}
\begin{csteps}
  \item Задать допустимую $x_0$, $\eps > 0$.
  \item Определить активный набор $I_k$, построить $A_k$ и $P_k$ (если
        активных нет — $P_k = E$).
  \item $g_k = \nabla f(x_k)$, $s_k = -P_k g_k$.
  \item $\norm{s_k} > \eps$ — направление найдено, к~шагу~5. Иначе вычислить
        множители Лагранжа: все $\lambda_i \ge 0$ — конец; иначе удалить
        ограничение с $\lambda_i < 0$ из $I_k$ и к~шагу~2.
  \item Найти $\alpha_{\max}$ и $\alpha_k$; $x_{k+1} = x_k + \alpha_k s_k$.
  \item Если шаг сделал новое ограничение активным — добавить его в $I_k$;
        $k := k+1$, к~шагу~2.
\end{csteps}

\paragraph{Ограничения.}
Нужна допустимая начальная точка; ограничения должны быть линейными (или
локально линеаризуемыми); при линейно зависимых строках активных ограничений
$A_kA_k\T$ вырождена; при $P_kg_k = 0$ нужна дополнительная проверка
множителей; для невыпуклой $f$ — локальный минимум.

\seefull{гл.~3, \S5--\S6 (численные примеры там же)}

\section{Билет 12. Метод условного градиента}

\ticketq{Метод условного градиента (метод Франка--Вульфа).}

\paragraph{Постановка.}
\[
  \varphi(x) \to \min_{x \in S} ,
  \qquad S \subset \R^n \ \text{— выпуклый компакт} ,
\]
$\varphi$ выпукла, непрерывно дифференцируема, $\nabla\varphi$ липшицев с
константой $R$. Метод относится к методам \textbf{линейной аппроксимации}:
направление ищется не проектированием градиента, а решением линейной задачи
на $S$.

\paragraph{Условие оптимальности для выпуклой задачи.}
\textbf{Теорема.} \textit{$x_* \in S$ — решение $\min_{x\in S}\varphi(x)$
$\iff$ $\nabla\T\varphi(x_*)(x - x_*) \ge 0$ $\forall x \in S$.}

\textsc{Док-во. Необходимость.} $S$ выпукло $\Rightarrow$
$x_* + \alpha(x - x_*) \in S$ при $\alpha \in [0,1]$. Функция
$\psi(\alpha) = \varphi\big(x_* + \alpha(x-x_*)\big)$ достигает минимума на
$[0,1]$ при $\alpha = 0$ $\Rightarrow$ $\psi'(0) \ge 0$, а
$\psi'(0) = \nabla\T\varphi(x_*)(x-x_*)$.

\textsc{Достаточность.} Для выпуклой дифференцируемой $\varphi$ верно
неравенство первого порядка
$\varphi(x) \ge \varphi(x_*) + \nabla\T\varphi(x_*)(x - x_*)$; второе
слагаемое неотрицательно $\Rightarrow$ $\varphi(x) \ge \varphi(x_*)$
$\forall x \in S$. $\blacksquare$

\paragraph{Вспомогательная задача и направление.}
В точке $x_k$ функция заменяется линейной аппроксимацией
$\varphi(x) \approx \varphi(x_k) + \nabla\T\varphi(x_k)(x - x_k)$; постоянные
слагаемые на точку минимума не влияют, поэтому решается
\[
  \nabla\T\varphi(x_k)\,x \to \min_{x \in S} ,
  \qquad y_k \ \text{— решение} ,
\]
\[
  s_k = y_k - x_k , \qquad
  \eta_k = \nabla\T\varphi(x_k)(y_k - x_k) ,
  \qquad
  x_{k+1} = x_k + \alpha_k s_k .
\]
Так как $y_k$ минимизирует линейную функцию, а $x_k \in S$, то
$\nabla\T\varphi(x_k)y_k \le \nabla\T\varphi(x_k)x_k$ $\Rightarrow$
$\eta_k \le 0$.

Если $S$ задано линейными ограничениями, вспомогательная задача — задача ЛП и
решается симплекс-методом; минимум линейной функции на выпуклом компакте
достигается на границе, а на многограннике — в вершине, поэтому $y_k$ всегда
«крайняя» и шаг $\alpha_k$ обязателен.

\paragraph{Теорема (направление спуска).}
\textit{$y_k \ne x_k$ и $\eta_k < 0$ $\Rightarrow$ $s_k = y_k - x_k$ —
направление спуска.}

\textsc{Док-во.} $\nabla\T\varphi(x_k)s_k = \eta_k < 0$ по определению.
$\blacksquare$

\paragraph{Теорема (критерий останова).}
\textit{Если $\eta_k = 0$, то $x_k$ — решение задачи.}

\textsc{Док-во.} $y_k$ — решение линейной задачи $\Rightarrow$
$\forall x \in S$:
$\nabla\T\varphi(x_k)y_k \le \nabla\T\varphi(x_k)x$. Вычитая
$\nabla\T\varphi(x_k)x_k$:
\[
  \underbrace{\nabla\T\varphi(x_k)(y_k - x_k)}_{=\,\eta_k\,=\,0}
  \ \le\ \nabla\T\varphi(x_k)(x - x_k)
  \ \Rightarrow\ \nabla\T\varphi(x_k)(x - x_k) \ge 0 \ \ \forall x \in S ,
\]
а это и есть условие оптимальности для выпуклой задачи. $\blacksquare$

\paragraph{Теорема (допустимость новой точки).}
\textit{$x_k, y_k \in S$, $S$ выпукло, $\alpha_k \in [0,1]$ $\Rightarrow$
$x_{k+1} \in S$.}

\textsc{Док-во.}
$x_{k+1} = x_k + \alpha_k(y_k - x_k) = (1-\alpha_k)x_k + \alpha_k y_k$ —
выпуклая комбинация точек $S$ $\Rightarrow$ принадлежит $S$. $\blacksquare$

Допустимость достаётся даром — никакого проектирования, в отличие от метода
Розена; достаточно держать $\alpha_k \in [0,1]$.

\paragraph{Алгоритм.}
\begin{csteps}
  \item Задать $x_0 \in S$, $\eps > 0$, $k = 0$.
  \item $g_k = \nabla\varphi(x_k)$.
  \item Решить ЛП: $y_k = \argmin_{x\in S} g_k\T x$; $s_k = y_k - x_k$,
        $\eta_k = g_k\T s_k$.
  \item $\eta_k \ge -\eps$ — останов, $x_k$ приближённо оптимальна.
  \item Выбрать $\alpha_k \in [0,1]$: дроблением
        $\alpha_k = \alpha_0\lambda^{\,j}$ с наименьшим $j$, при котором
        $\varphi(x_k + \alpha_k s_k) - \varphi(x_k) \le
         \tfrac12\alpha_k\eta_k$; либо точной минимизацией
        $\argmin_{0\le\alpha\le1}\varphi(x_k+\alpha s_k)$; либо по заданной
        последовательности $\alpha_k = \frac{2}{k+2}$.
  \item $x_{k+1} = x_k + \alpha_k s_k$, $k := k+1$, к~шагу~2.
\end{csteps}

\paragraph{Теорема о сходимости и скорость $O(1/k)$.}
\textit{$\varphi(x_k) \to \varphi(x_*) = \min_S\varphi$, причём
$\varphi(x_k) - \varphi(x_*) \le D/k$, $D > 0$.}

\textsc{Док-во.} Пусть $d = \operatorname{diam}S$,
$G = \max_S\norm{\nabla\varphi}$, тогда $\abs{\eta_k} \le Gd$. По лемме о
дроблении шаг отделён от нуля,
$\alpha_k \ge \lambda\min\{\alpha_0,\ \abs{\eta_k}/(2Rd^2)\}$, поэтому
\[
  \varphi(x_k) - \varphi(x_{k+1}) \ \ge\ \tfrac12\alpha_k\abs{\eta_k}
   \ \ge\ c\,\eta_k^2 ,
  \qquad
  c = \tfrac{\lambda}{2}\min\Big\{\tfrac{\alpha_0}{Gd},\ \tfrac{1}{2Rd^2}\Big\} .
\]
$\varphi(x_k)$ убывает и ограничена снизу $\Rightarrow$
$\sum_k\eta_k^2 < \infty$ $\Rightarrow$ $\eta_k \to 0$. По выпуклости и
определению $\eta_k$ как минимума:
\[
  \varphi(x_k) - \varphi(x_*) \le \nabla\T\varphi(x_k)(x_k - x_*)
   \le -\eta_k = \abs{\eta_k}
   \ \Rightarrow\ \varphi(x_k) \to \varphi(x_*) .
\]
Скорость: обозначив $\delta_k = \varphi(x_k) - \varphi_*$, имеем
$\delta_k \le \abs{\eta_k}$ и $\delta_k - \delta_{k+1} \ge c\eta_k^2
\ge c\delta_k^2$, откуда
\[
  \frac{1}{\delta_{k+1}} - \frac{1}{\delta_k}
   = \frac{\delta_k - \delta_{k+1}}{\delta_k\delta_{k+1}}
   \ \ge\ \frac{c\delta_k^2}{\delta_k\delta_{k+1}} \ \ge\ c
  \quad (\delta_{k+1} \le \delta_k)
  \ \Rightarrow\ \delta_k \le \frac{1}{ck + 1/\delta_0} \le \frac{D}{k} .
  \ \blacksquare
\]
Скорость \textbf{сублинейная} — заметно хуже геометрической сходимости методов
безусловной минимизации. Плата за это — крайняя дешевизна итерации: одна
линейная задача вместо проектирования.

\paragraph{Ограничения.}
На каждой итерации решается задача ЛП на $S$; сходимость медленная около
решения; при неограниченном $S$ линейная подзадача может не иметь конечного
решения; для невыпуклой $\varphi$ — только стационарная точка.

\seefull{гл.~3, \S8 (численный пример там же)}

\section{Билет 13. Метод Ньютона решения задач с ограничениями}

\ticketq{Метод Ньютона решения задач с ограничениями.}

\paragraph{Постановка.}
\[
  \varphi(x) \to \min , \qquad x \in S ,
\]
$\varphi$ выпукла и дважды дифференцируема, $S$ — выпуклый компакт,
$H(x) = \nabla^2\varphi(x)$. Итерации $x_{k+1} = x_k + \alpha_k s_k$.

\textbf{Идея.} Это метод Ньютона, перенесённый на задачу с ограничениями: там,
где метод условного градиента (билет~12) минимизирует на $S$ \emph{линейную}
модель, здесь минимизируется \emph{квадратичная}.

\paragraph{Квадратичная модель и направление.}
\[
  \psi_k(x) = \nabla\T\varphi(x_k)(x - x_k)
            + \tfrac12 (x - x_k)\T H(x_k)(x - x_k) ,
\]
\[
  y_k = \argmin_{x \in S}\psi_k(x) , \qquad
  \eta_k = \psi_k(y_k) \le \psi_k(x_k) = 0 , \qquad
  s_k = y_k - x_k .
\]

\paragraph{Лемма (вспомогательные факты).}
\textit{1) $x_k + \alpha s_k \in S$ при $\alpha \in [0,1]$;
2) $\nabla\T\varphi(x_k)s_k \le \eta_k \le 0$;
3) при сильной выпуклости с константой $m$: $\eta_k \le -\tfrac{m}{2}\norm{s_k}^2$.}

\textsc{Док-во.}
1) $x_k + \alpha s_k = (1-\alpha)x_k + \alpha y_k$ — выпуклая комбинация точек
$S$.

2) $\eta_k = \nabla\T\varphi(x_k)s_k + \tfrac12 s_k\T H(x_k)s_k$, второе
слагаемое $\ge 0$ ($H \succeq 0$ по выпуклости) $\Rightarrow$
$\nabla\T\varphi(x_k)s_k \le \eta_k$; а $\eta_k \le \psi_k(x_k) = 0$.

3) $\psi_k$ сильно выпукла с той же $m$, $y_k$ — её минимум на $S$
$\Rightarrow$ $\nabla\T\psi_k(y_k)(x_k - y_k) \ge 0$, и
\[
  0 = \psi_k(x_k) \ \ge\ \psi_k(y_k)
    + \underbrace{\nabla\T\psi_k(y_k)(x_k-y_k)}_{\ge\,0}
    + \tfrac{m}{2}\norm{x_k - y_k}^2
  \ \ge\ \eta_k + \tfrac{m}{2}\norm{s_k}^2 . \ \blacksquare
\]

\paragraph{Выбор шага.}
$\alpha_0 = 1$, $\eps, \lambda \in (0,1)$; $\alpha_k = \alpha_0\lambda^{i_k}$ —
первое значение, удовлетворяющее условию дробления
\[
  \varphi(x_k + \alpha s_k) - \varphi(x_k) \ \le\ \alpha\,\eps\,\eta_k
\]
(либо наискорейший спуск вдоль $s_k$: $\min_{\alpha>0}\varphi(x_k+\alpha s_k)$).

\textbf{Лемма (дробление конечно).} $H$ непрерывен на компакте $\Rightarrow$
$\norm{H(x)} \le M$, откуда
\[
  \varphi(x_k + \alpha s_k) \le \varphi(x_k) + \alpha\eta_k
    + \tfrac{M}{2}\alpha^2\norm{s_k}^2 ,
\]
и условие выполнено, как только
\[
  \alpha\eta_k + \tfrac{M}{2}\alpha^2\norm{s_k}^2 \le \alpha\eps\eta_k
  \iff
  \alpha \le -\frac{2(1-\eps)\eta_k}{M\norm{s_k}^2} > 0 \quad (\eta_k < 0) .
\]
Значит дробление конечно, и принятый шаг отделён от нуля:
$\alpha_k \ge \min\{1,\ \tfrac{2\lambda(1-\eps)}{MC^2}\abs{\eta_k}\}$,
$C = \operatorname{diam}S$. $\blacksquare$

\paragraph{Критерий останова.}
$\eta_k = 0$ $\Rightarrow$ $x_k$ минимизирует $\psi_k$ на $S$, и условие
оптимальности $\nabla\T\psi_k(x_k)(x - x_k) = \nabla\T\varphi(x_k)(x-x_k) \ge 0$
$\forall x \in S$ для выпуклой $\varphi$ означает, что $x_k$ — решение.

\paragraph{Теоремы о сходимости.}
\textbf{Т1.} \textit{$\varphi(x_k)$ монотонно убывает и
$\lim_k \varphi(x_k) = \varphi(x_*) = \min_{x\in S}\varphi(x)$ при любом $x_0$.}

\textsc{Схема.} Из условия дробления
$\varphi(x_k) - \varphi(x_{k+1}) \ge \eps\alpha_k\abs{\eta_k} \ge 0$:
последовательность убывает, ограничена снизу ($S$ компакт) $\Rightarrow$
сходится $\Rightarrow$ $\alpha_k\abs{\eta_k} \to 0$, а с оценкой шага снизу
$\eta_k \to 0$. Далее для $x \in S$, $t \in [0,1]$ из $\psi_k(x_k+t(x-x_k)) \ge \eta_k$
следует $t\nabla\T\varphi(x_k)(x-x_k) \ge \eta_k - \tfrac{t^2}{2}MC^2$; при
$t = \sqrt{\abs{\eta_k}}$ получаем
$\nabla\T\varphi(x_k)(x_* - x_k) \ge -\sqrt{\abs{\eta_k}}\big(1+\tfrac{MC^2}{2}\big)$,
и по выпуклости
$\varphi(x_k) - \varphi(x_*) \le \sqrt{\abs{\eta_k}}\big(1+\tfrac{MC^2}{2}\big) \to 0$.
$\blacksquare$

\textbf{Т2.} \textit{Если $m\norm{y}^2 \le y\T H(x)y \le M\norm{y}^2$,
$0 < m \le M$, то $x_k \to x_*$ \textbf{сверхлинейно}.}

\textbf{Т3.} \textit{Если дополнительно $\norm{H(x)-H(y)} \le R\norm{x-y}$, то
скорость \textbf{квадратичная}.}

\textsc{Схема Т2--Т3.} Вблизи $x_*$ при $\eps < \tfrac12$ принимается полный шаг
$\alpha_k = 1$, то есть $x_{k+1} = y_k$. Складывая условия оптимальности
$\big(\nabla\varphi(x_k) + H_k(y_k - x_k)\big)\T(x_* - y_k) \ge 0$ и
$\nabla\T\varphi(x_*)(y_k - x_*) \ge 0$ и обозначая
$r_k = \nabla\varphi(x_k) - \nabla\varphi(x_*) - H_k(x_k - x_*)$:
\[
  m\norm{y_k - x_*}^2 \le \norm{r_k}\,\norm{y_k - x_*}
  \ \Rightarrow\ \norm{x_{k+1} - x_*} \le \frac{\norm{r_k}}{m} .
\]
Без липшицевости $H$: $\norm{r_k} = o(\norm{x_k - x_*})$ — сверхлинейная
сходимость; при $\norm{H(x)-H(y)} \le R\norm{x-y}$:
$\norm{r_k} \le \tfrac{R}{2}\norm{x_k-x_*}^2$ — квадратичная. $\blacksquare$

\paragraph{Алгоритм.}
\begin{csteps}
  \item Задать $x_0 \in S$, $\eps, \lambda \in (0,1)$, $k = 0$.
  \item Построить $\psi_k(x)$ по $\nabla\varphi(x_k)$ и $H(x_k)$.
  \item Решить $y_k = \argmin_{x\in S}\psi_k(x)$; найти
        $\eta_k = \psi_k(y_k)$, $s_k = y_k - x_k$.
  \item $\eta_k = 0$ — конец, $x_k$ оптимальна.
  \item Найти $\alpha_k$ дроблением от $\alpha_0 = 1$ (либо наискорейшим
        спуском вдоль $s_k$).
  \item $x_{k+1} = x_k + \alpha_k s_k$, $k := k+1$, к~шагу~2.
\end{csteps}

\textbf{Замечание.} Для задачи с ограничениями-\emph{равенствами} $h(x) = 0$
тот же метод Ньютона записывают иначе: условия Лагранжа
$\nabla f + J\T\lambda = 0$, $h = 0$ рассматривают как систему $F(x,\lambda) = 0$
и применяют к ней метод Ньютона, получая на каждом шаге блочную систему ККТ
$\big(\begin{smallmatrix} W_k & J_k\T \\ J_k & 0\end{smallmatrix}\big)
\big(\begin{smallmatrix} p_k \\ \Delta\lambda_k\end{smallmatrix}\big)
= -\big(\begin{smallmatrix} \nabla_x L \\ h \end{smallmatrix}\big)$.

\seefull{гл.~3, \S10 (а вариант с системой ККТ — \S13)}

\section{Билет 14. Метод отсекающей гиперплоскости}

\ticketq{Метод отсекающей гиперплоскости.}

\paragraph{Постановка.}
\[
  \varphi_0(x) = c\T x \to \min , \qquad x \in \Omega ,
  \qquad \Omega = \{\, x \in \R^n \mid \varphi(x) \le 0 \,\} ,
\]
$\varphi$ выпукла, $\Omega \ne \varnothing$ — компакт. \textbf{Идея:}
заменить сложное $\Omega$ простым объемлющим многогранником и последовательно
добавлять линейные ограничения, отсекающие недопустимые точки.

\textbf{Приведение к этому виду.} Система $\varphi_i(x) \le 0$,
$i = \overline{1,m}$, сворачивается в одно ограничение
$\varphi(x) = \max_i \varphi_i(x)$ (максимум конечного числа выпуклых функций
выпуклен). Нелинейная выпуклая цель заменяется линейной вводом переменной:
$x^{(n+1)} \to \min$ при $\varphi_0(x) - x^{(n+1)} \le 0$.

\paragraph{Отсекающая гиперплоскость.}
Гиперплоскость $a\T x = \beta$ называется \textbf{отсекающей} для $y \notin \Omega$,
если
\[
  a\T x \le \beta \ \ \forall x \in \Omega
  \qquad\text{и}\qquad
  a\T y > \beta ,
\]
то есть новое ограничение сохраняет всё допустимое множество и выбрасывает $y$.

\textbf{Субградиент.} Вектор $p$ — субградиент выпуклой $\varphi$ в точке $y$,
если $\varphi(x) \ge \varphi(y) + p\T(x-y)$ $\forall x$. Для
дифференцируемой $\varphi$ субградиент совпадает с $\nabla\varphi(y)$.

\paragraph{Теорема (построение отсечения).}
\textit{Пусть $\varphi$ выпукла, $p \in \partial\varphi(y)$ и
$\varphi(y) > 0$. Тогда}
\[
  \varphi(y) + p\T(x - y) \le 0
\]
\textit{— отсекающее неравенство для точки $y$.}

\textsc{Док-во.} \emph{Допустимые точки сохраняются:} для $x \in \Omega$ по
определению субградиента
\[
  \varphi(y) + p\T(x-y) \ \le\ \varphi(x) \ \le\ 0 .
\]
\emph{Точка $y$ отсекается:} подстановка $x = y$ даёт
$\varphi(y) + p\T(y-y) = \varphi(y) > 0$, то есть неравенство нарушено.
$\blacksquare$

\paragraph{Внешние приближения.}
Строится последовательность многогранников
$S_0 \supset S_1 \supset \dots \supset \Omega$. На $k$-м шаге решается
линейная задача на $S_k$, её решение $x_k$; если $\varphi(x_k) > 0$, строится
субградиент $a_{k+1} \in \partial\varphi(x_k)$ и, при
$b^{(k+1)} = -\varphi(x_k) + a_{k+1}\T x_k$,
\[
  \varphi(x) \ \ge\ a_{k+1}\T x - b^{(k+1)} ,
  \qquad
  S_{k+1} = S_k \cap \big\{\, a_{k+1}\T x - b^{(k+1)} \le 0 \,\big\} .
\]

\textbf{Лемма (вложенность).} \textit{$\Omega \subset S_k$ и отсечение
корректно $\Rightarrow$ $\Omega \subset S_{k+1} \subset S_k$.}

\textsc{Док-во.} $x \in \Omega$ $\Rightarrow$ $x \in S_k$ (по условию) и
$a_{k+1}\T x - b^{(k+1)} \le \varphi(x) \le 0$ $\Rightarrow$ $x$ лежит в
пересечении, то есть в $S_{k+1}$. Текущая же $x_k$ выбрасывается:
$a_{k+1}\T x_k - b^{(k+1)} = \varphi(x_k) > 0$. $\blacksquare$

\paragraph{Алгоритм.}
\begin{csteps}
  \item Построить объемлющий многогранник-компакт $S_0 \supset \Omega$,
        заданный линейными ограничениями; $k = 0$.
  \item Решить $c\T x \to \min$, $x \in S_k$; получить $x_k$.
  \item $\varphi(x_k) \le 0$ — точка допустима, она и есть решение, конец.
  \item Иначе взять $a_{k+1} \in \partial\varphi(x_k)$, вычислить
        $b^{(k+1)} = -\varphi(x_k) + a_{k+1}\T x_k$ и положить
        $S_{k+1} = S_k \cap \{a_{k+1}\T x \le b^{(k+1)}\}$.
  \item $k := k+1$; останов при $\norm{x_{k+1} - x_k} < \eps$, иначе
        к~шагу~2.
\end{csteps}

\paragraph{Особенности.}
\textbf{Линейная подзадача.} От шага к шагу меняется только \emph{число}
ограничений $\Rightarrow$ применим симплекс-метод.

\textbf{Двойственность.} Добавление ограничения к прямой задаче — это
добавление одной \emph{переменной} к двойственной. Поэтому предыдущее
двойственное решение, дополненное нулевой компонентой, остаётся допустимым и
служит готовым начальным приближением: симплекс-метод продолжает работу, а не
стартует заново. В этом главный вычислительный выигрыш.

\paragraph{Теорема о сходимости.}
\textit{Пусть $\varphi$ непрерывна и выпукла, $\Omega$ и $S_0$ компактны,
субградиенты ограничены: $\norm{a_k} \le R$. Тогда любая предельная точка
$\{x_k\}$ — оптимальное решение, причём $\varphi(x_k) \to 0$.}

\textsc{Док-во.} $S_{k+1} \subseteq S_k \subseteq S_0$ и $S_0$ компактно
$\Rightarrow$ предельные точки существуют.

\emph{1)} Если при каком-то $k$ оказалось $\varphi(x_k) \le 0$, то
$x_k \in \Omega \subseteq S_k$ и $\varphi_0(x_k) = \min_{S_k}\varphi_0 \le
\varphi_0(x)$ $\forall x \in \Omega$ — точка уже оптимальна.

\emph{2)} Пусть $\varphi(x_k) > 0$ всегда и $x_k \to \bar{x}$ по подмножеству
$K$. Для $j < k$ имеем $x_k \in S_{j+1}$, то есть
$a_{j+1}\T(x_k - x_j) + \varphi(x_j) \le 0$, откуда
\[
  \varphi(x_j) \ \le\ \norm{a_{j+1}}\,\norm{x_j - x_k}
   \ \le\ R\,\norm{x_j - x_k} .
\]
Для $j,k \in K$ расстояние $\norm{x_j - x_k} \to 0$ $\Rightarrow$
$\varphi(x_j) \to 0$, и по непрерывности $\varphi(\bar{x}) = 0$, то есть
$\bar{x} \in \Omega$.

\emph{3)} $\varphi_0(x_k) = \min_{S_k}\varphi_0 \le \varphi_0(x)$
$\forall x \in \Omega$; переходя к пределу,
$\varphi_0(\bar{x}) \le \varphi_0(x)$ — точка $\bar{x}$ оптимальна.
$\blacksquare$

Содержательно: все $x_k$ недопустимы, поэтому дают \emph{нижние} оценки
оптимума, а предельная точка допустима и даёт верхнюю — метод зажимает оптимум
с двух сторон.

\paragraph{Ограничения.}
Может потребоваться много отсечений; подзадача растёт с каждой итерацией; при
слабых отсечениях область сокращается медленно; требуется выпуклость — для
невыпуклого $\Omega$ касательная гиперплоскость может отсечь часть допустимого
множества; при недифференцируемой $\varphi$ нужно уметь строить субградиенты.

\seefull{гл.~3, \S2 и \S7 (численный пример в \S7)}

\section{Билет 15. Метод штрафных функций}

\ticketq{Метод штрафных функций.}

\paragraph{Постановка.}
\[
  \varphi(x) \to \min , \qquad
  x \in S = \{\, x \mid \psi_i(x) \le 0,\ i = \overline{1,m};\
    \chi_j(x) = 0,\ j = \overline{1,p} \,\} .
\]
Метод \textbf{внешних} штрафов сводит задачу к последовательности задач
\emph{безусловной} минимизации: нарушение ограничений не запрещается, а
штрафуется. Точки при этом могут быть недопустимыми.

\paragraph{Штрафная функция.}
$\alpha(x) \ge 0$, причём $\alpha(x) = 0 \iff x \in S$. Стандартные формы
($p > 1$, обычно $p = 2$):
\[
  \text{одно неравенство } \psi \le 0 :\quad
  \alpha(x) = \begin{cases} 0, & \psi(x) \le 0 , \\
                            \psi(x)^p, & \psi(x) > 0 ; \end{cases}
\]
\[
  \text{одно равенство } \chi = 0 :\quad \alpha(x) = \abs{\chi(x)}^p ;
\]
\[
  \text{общий случай:}\quad
  \alpha(x) = \sum_{i=1}^{m}\xi_i \max\{0, \psi_i(x)\}^p
            + \sum_{j=1}^{p}\eta_j \abs{\chi_j(x)}^p ,
  \qquad \xi_i, \eta_j > 0 .
\]

\paragraph{Вспомогательная задача.}
\[
  \varphi(x) + \mu\,\alpha(x) \to \min_{x \in V} , \quad \mu > 0 ,
  \qquad x_\mu \ \text{— решение} ,
  \qquad P(\mu) = \varphi(x_\mu) + \mu\,\alpha(x_\mu) .
\]

\paragraph{Лемма (свойства вспомогательных задач).}
\textit{Пусть все функции непрерывны и при каждом $\mu > 0$ решение $x_\mu$
существует. Тогда}
\begin{csteps}
  \item $\min\limits_{x\in S}\varphi(x) \ \ge\ \sup\limits_{\mu>0}P(\mu)$;
  \item $\alpha(x_\mu)$ — невозрастающая функция $\mu$;
  \item $\varphi(x_\mu)$ — неубывающая функция $\mu$;
  \item $P(\mu)$ — неубывающая функция $\mu$;
  \item $\lim\limits_{\mu\to\infty}P(\mu) = \sup\limits_{\mu>0}P(\mu)$.
\end{csteps}

\textsc{Док-во.}
\textbf{1)} Для $x \in S$ штраф нулевой, а $x_\mu$ минимизирует
$\varphi + \mu\alpha$ на $V \supset S$:
\[
  P(\mu) = \min_{z\in V}\big(\varphi(z)+\mu\alpha(z)\big)
   \le \varphi(x) + \mu\underbrace{\alpha(x)}_{=\,0} = \varphi(x) ;
\]
минимум по $x \in S$ и супремум по $\mu$ дают~1).

\textbf{2)} Пусть $0 < \lambda < \mu$. Из минимальности $x_\lambda$ и $x_\mu$:
\[
  \varphi(x_\lambda) + \lambda\alpha(x_\lambda)
    \le \varphi(x_\mu) + \lambda\alpha(x_\mu) , \qquad
  \varphi(x_\mu) + \mu\alpha(x_\mu)
    \le \varphi(x_\lambda) + \mu\alpha(x_\lambda) .
\]
Складывая: $(\mu-\lambda)\big(\alpha(x_\mu) - \alpha(x_\lambda)\big) \le 0$
$\Rightarrow$ $\alpha(x_\mu) \le \alpha(x_\lambda)$.

\textbf{3)} Из первого неравенства:
$\varphi(x_\lambda) - \varphi(x_\mu) \le
 \lambda\big(\alpha(x_\mu)-\alpha(x_\lambda)\big) \le 0$.

\textbf{4)}
$P(\lambda) \le \varphi(x_\mu) + \lambda\alpha(x_\mu)
 \le \varphi(x_\mu) + \mu\alpha(x_\mu) = P(\mu)$.

\textbf{5)} $P$ монотонна и ограничена сверху величиной $\min_S\varphi$ (п.~1)
$\Rightarrow$ предел равен супремуму. $\blacksquare$

\paragraph{Теорема о сходимости.}
\textit{Пусть решение $x_*$ существует, $\mu_k \to \infty$ и $\{x_{\mu_k}\}$
лежит в компакте. Тогда}
\[
  \min_{S}\varphi = \lim_k \varphi(x_{\mu_k}) , \qquad
  x_{\mu_k} \to x_* , \qquad
  \mu_k\,\alpha(x_{\mu_k}) \to 0 .
\]

\textsc{Док-во.} По лемме $\varphi(x_\mu) \le P(\mu) \le \varphi_*$, где
$\varphi_* = \min_S\varphi$, поэтому
\[
  \mu\,\alpha(x_\mu) = P(\mu) - \varphi(x_\mu)
   \le \varphi_* - \varphi(x_{\mu_1})
\]
ограничено сверху константой $\Rightarrow$ $\alpha(x_\mu) \to 0$ при
$\mu \to \infty$. Для предельной точки $\bar{x}$ по непрерывности
$\alpha(\bar{x}) = 0$ $\Rightarrow$ $\bar{x} \in S$ $\Rightarrow$
$\varphi(\bar{x}) \ge \varphi_*$. С другой стороны
$\varphi(\bar{x}) = \lim\varphi(x_{\mu_k}) \le \varphi_*$. Значит
$\varphi(\bar{x}) = \varphi_*$ — точка $\bar{x}$ оптимальна. Отсюда же
$P(\mu) \to \varphi_*$ и $\mu\alpha(x_\mu) = P(\mu) - \varphi(x_\mu) \to 0$.
$\blacksquare$

\paragraph{Алгоритм.}
\begin{csteps}
  \item Задать $\mu_1 > 0$, $\eps > 0$, $\lambda > 1$, начальную $x_1$ (может
        быть недопустимой), $k = 1$.
  \item Решить $\min_x\big[\varphi(x) + \mu_k\alpha(x)\big]$ любым методом
        безусловной минимизации (билеты~6--9); получить $x_{\mu_k}$.
  \item Если $\mu_k\alpha(x_{\mu_k}) < \eps$ — конец.
  \item Иначе $\mu_{k+1} = \lambda\mu_k$, $k := k+1$, к~шагу~2.
\end{csteps}

\textbf{Про критерий останова.} $\mu\alpha(x_\mu) < \eps$ означает
$\varphi_* - \varphi(x_\mu) \le \varphi_* - P(\mu) + \mu\alpha(x_\mu)$, где
$\varphi_* - P(\mu) \to 0$; практически это малое нарушение ограничений.

\paragraph{Внешние и внутренние штрафы.}
Здесь $\mu \to +\infty$ и точки могут быть недопустимыми. Во внутреннем
(барьерном, билет~16) методе наоборот: точки обязаны лежать строго внутри, а
параметр $\mu \to +0$.

\paragraph{Ограничения.}
При больших $\mu$ вспомогательная функция плохо обусловлена и минимизируется
тяжелее; её минимум лежит вне $S$; точное решение требует $\mu \to \infty$,
что недостижимо; слишком быстрый рост $\mu$ делает счёт неустойчивым,
слишком медленный — долгим.

\seefull{гл.~3, \S3 и \S11 (численный пример в \S11)}

\section{Билет 16. Метод барьерных функций}

\ticketq{Метод барьерных функций.}

\paragraph{Постановка.}
\[
  \varphi(x) \to \min , \qquad
  S = \{\, x \mid \psi_i(x) \le 0,\ i = \overline{1,m} \,\} ,
  \qquad
  T = \{\, x \mid \psi_i(x) < 0 \,\} \ne \varnothing
\]
($T$ — внутренность $S$). Метод внутренних штрафов: итерации всё время
остаются строго внутри $S$, а к цели добавляется барьер, растущий до
бесконечности у границы.

\textbf{Важно:} метод неприменим к ограничениям-\emph{равенствам} — у
множества $\chi(x) = 0$ внутренности нет вовсе. Равенства обрабатывают внешним
штрафом (билет~15), барьером — только неравенства.

\paragraph{Барьерная функция.}
$B(x) \to +\infty$ при $x \to \partial S$ изнутри. Два стандартных вида
($\xi_i > 0$):
\[
  B(x) = -\sum_{i=1}^{m}\xi_i\,\frac{1}{\psi_i(x)} ,
  \qquad
  B(x) = -\sum_{i=1}^{m}\xi_i\,\ln\big(-\psi_i(x)\big) .
\]
Обратный барьер неудобен для дифференцирования, логарифмический удобнее.

\textbf{Почему барьер работает.} $\psi_i(x) \to 0-$ $\Rightarrow$
$-\psi_i(x) \to 0+$ $\Rightarrow$ $\ln(-\psi_i) \to -\infty$ $\Rightarrow$
$-\ln(-\psi_i) \to +\infty$ (аналогично $-1/\psi_i \to +\infty$). Кроме того,
$B$ определена \emph{только} на $T$, поэтому всякая точка с конечным
значением $\varphi + \mu B$ автоматически строго допустима — проверять
допустимость отдельно не нужно.

\paragraph{Вспомогательная задача.}
\[
  \varphi(x) + \mu B(x) \to \min_{x \in T} , \quad \mu > 0 ,
  \qquad x_\mu \ \text{— решение} ,
  \qquad \Theta(\mu) = \varphi(x_\mu) + \mu B(x_\mu) ,
\]
причём $\mu \to +0$ (в отличие от внешнего штрафа, где $\mu \to +\infty$).
Точки $x_\mu$ образуют \textbf{центральную траекторию}.

\paragraph{Лемма (свойства вспомогательных задач).}
\textit{Считаем $B \ge 0$ на $T$. Тогда}
\begin{csteps}
  \item $\min\limits_{S}\varphi \ \le\ \Theta(\mu)$ при всех $\mu > 0$;
  \item $B(x_\mu)$ не убывает при $\mu \to 0$
        ($\mu_1 > \mu_2 \Rightarrow B(x_{\mu_1}) \le B(x_{\mu_2})$);
  \item $\varphi(x_\mu)$ не возрастает при $\mu \to 0$;
  \item $\Theta(\mu)$ — неубывающая функция $\mu$;
  \item $\inf\limits_{\mu>0}\Theta(\mu) = \lim\limits_{\mu\to0}\Theta(\mu)$.
\end{csteps}

\textsc{Док-во.} Пусть $\mu_1 > \mu_2 > 0$, $x_i = x_{\mu_i}$.

\textbf{1)} $x_\mu \in T \subset S$ и $B \ge 0$ $\Rightarrow$
$\Theta(\mu) \ge \varphi(x_\mu) \ge \min_S\varphi$.

\textbf{2)} Из минимальности $x_1$ и $x_2$:
\[
  \varphi(x_1) + \mu_1 B(x_1) \le \varphi(x_2) + \mu_1 B(x_2) , \qquad
  \varphi(x_2) + \mu_2 B(x_2) \le \varphi(x_1) + \mu_2 B(x_1) .
\]
Складывая: $(\mu_1 - \mu_2)\big(B(x_1) - B(x_2)\big) \le 0$ $\Rightarrow$
$B(x_1) \le B(x_2)$.

\textbf{3)} Из второго неравенства:
$\varphi(x_2) - \varphi(x_1) \le \mu_2\big(B(x_1)-B(x_2)\big) \le 0$.

\textbf{4)}
$\Theta(\mu_2) \le \varphi(x_1) + \mu_2 B(x_1)
 \le \varphi(x_1) + \mu_1 B(x_1) = \Theta(\mu_1)$.

\textbf{5)} $\Theta$ монотонна и ограничена снизу (п.~1) $\Rightarrow$
нижняя грань достигается в пределе. $\blacksquare$

\paragraph{Теорема о сходимости.}
\textit{Пусть $T \ne \varnothing$, функции непрерывны на $T$, при каждом
$\mu > 0$ решение $x_\mu$ существует, $\{x_\mu\}$ лежит в компакте и всякая
точка $S$ приближается точками $T$ (то есть $\forall x \in S$,
$\forall\eps > 0$ $\exists y \in T$: $\norm{y - x} < \eps$). Тогда предельные
точки $\{x_\mu\}$ — решения задачи, причём}
\[
  \mu B(x_\mu) \to 0 , \qquad
  \min_S \varphi = \lim_{\mu\to0}\Theta(\mu) .
\]

\textsc{Док-во.} По условию для любого $\eps$ найдётся $y \in T$ с
$\varphi(y) < \varphi_* + \eps$, где $\varphi_* = \min_S\varphi$. Тогда
\[
  \Theta(\mu) \le \varphi(y) + \mu B(y) \xrightarrow[\mu\to0]{} \varphi(y)
  \ \Rightarrow\ \lim_{\mu\to0}\Theta(\mu) \le \varphi_* + \eps
  \quad \forall\eps
  \ \Rightarrow\ \lim\Theta \le \varphi_* ,
\]
а по п.~1 леммы $\lim\Theta \ge \varphi_*$. Значит
$\lim_{\mu\to0}\Theta(\mu) = \varphi_*$. Далее
\[
  0 \le \mu B(x_\mu) = \Theta(\mu) - \varphi(x_\mu)
    \le \Theta(\mu) - \varphi_* \to 0 .
\]
Для предельной точки $\bar{x}$: $\bar{x} \in S$ (замкнутость) и
$\varphi(\bar{x}) = \lim\varphi(x_\mu) \le \lim\Theta = \varphi_*$, то есть
$\bar{x}$ — решение. $\blacksquare$

\paragraph{Алгоритм.}
\begin{csteps}
  \item Задать $x_1 \in T$ (строго допустимую), $\mu_1 > 0$, $\eps > 0$,
        $\beta \in (0,1)$, $k = 1$.
  \item Найти $x_{\mu_k} = \argmin_{x\in T}\big[\varphi(x) + \mu_k B(x)\big]$,
        стартуя с предыдущей точки.
  \item Если $\mu_k B(x_{\mu_k}) < \eps$ — конец.
  \item Иначе $\mu_{k+1} = \beta\mu_k$, $k := k+1$, к~шагу~2.
\end{csteps}

\paragraph{Ограничения.}
Нужна строго допустимая $x_1$; метод неприменим при $T = \varnothing$; точки
не могут лежать на границе, где барьер не определён, — а оптимум, как правило,
именно там; при малых $\mu$ функция $\varphi + \mu B$ плохо обусловлена у
границы; равенства напрямую не поддерживаются.

\seefull{гл.~3, \S3 и \S12 (численный пример в \S12)}

\section{Билет 17. Метод приведённого градиента Вульфа}

\ticketq{Метод приведённого градиента Вульфа.}

\paragraph{Постановка.}
\[
  \varphi(x) \to \min , \qquad Ax = b , \qquad x \ge 0 ,
\]
$\varphi$ выпукла и дифференцируема, $A$ размера $m\times n$, $m < n$,
$\rank A = m$.

\textbf{Условие невырожденности:} любые $m$ столбцов $A$ линейно независимы и
любая крайняя точка имеет ровно $m$ положительных компонент.

\textbf{Идея:} равенства используются для выражения базисных переменных через
небазисные; градиент пересчитывается с учётом этой зависимости; движение идёт
по антиприведённому градиенту.

\paragraph{Разбиение переменных.}
$x = (x_B, x_N)$, $A = (B,N)$, $B$ невырождена, $x_B > 0$:
\[
  Ax = b \iff Bx_B + Nx_N = b \iff x_B = B^{-1}(b - Nx_N) .
\]

\textbf{Лемма (сохранение равенств).} $Ax = b$, $Ad = 0$ $\Rightarrow$
$A(x+\alpha d) = Ax + \alpha Ad = b$. $\blacksquare$

Разбивая $d = (d_B, d_N)$, из $Ad = 0$ получаем $Bd_B + Nd_N = 0$, то есть
\[
  d_B = -B^{-1}N\,d_N :
\]
достаточно выбрать небазисную часть, базисная определяется однозначно.

\paragraph{Приведённый градиент.}
\[
  n\T = \nabla\T\varphi - \nabla_B\T\varphi\,B^{-1}A
      = \big(0,\ \nabla_N\T\varphi - \nabla_B\T\varphi\,B^{-1}N\big) ,
\]
то есть $n_N = \nabla_N\varphi - N\T B^{-\mathsf{T}}\nabla_B\varphi$;
компоненты обозначаем $n^{(j)}$, $j \in N$.

\textbf{Теорема.} \textit{Для допустимого направления ($Ad = 0$):}
$\nabla\T\varphi\,d = n_N\T d_N$.

\textsc{Док-во.} Подставим $d_B = -B^{-1}Nd_N$:
\[
  \nabla\T\varphi\,d
   = \nabla_B\T\varphi\,d_B + \nabla_N\T\varphi\,d_N
   = \big(\nabla_N\T\varphi - \nabla_B\T\varphi\,B^{-1}N\big)d_N
   = n_N\T d_N . \ \blacksquare
\]
Задача свелась к движению по небазисным переменным, и роль градиента в этом
пространстве играет $n_N$.

\paragraph{Направление Вульфа.}
Двигаться просто по $-n_N$ нельзя: при $x_j = 0$ отрицательная компонента
нарушила бы $x_j \ge 0$. Поэтому
\[
  d_j =
  \begin{cases}
    -\,n^{(j)} , & n^{(j)} \le 0 , \\[2pt]
    -\,x^{(j)}\,n^{(j)} , & n^{(j)} > 0 ,
  \end{cases}
  \qquad j \in N ,
  \qquad d_B = -B^{-1}N\,d_N .
\]
Множитель $x^{(j)}$ во втором случае и решает проблему: при $x_j = 0$
компонента $d_j$ обнуляется сама.

\textbf{Допустимость.} Для $j \in N$ с $x_j = 0$: либо $n^{(j)} \le 0$ и
$d_j = -n^{(j)} \ge 0$, либо $n^{(j)} > 0$ и $d_j = -x_jn^{(j)} = 0$ — в
обоих случаях $d_j \ge 0$. Базисные компоненты строго положительны
(невырожденность), поэтому при малом $\alpha$ имеем $x + \alpha d \ge 0$.

\paragraph{Теорема (направление спуска).}
\textit{$\nabla\T\varphi\,d \le 0$, причём равенство $\iff d_N = 0$.}

\textsc{Док-во.} По предыдущей теореме и правилу Вульфа:
\[
  n_N\T d_N
   = -\sum_{n^{(j)} \le 0}\big(n^{(j)}\big)^2
     - \sum_{n^{(j)} > 0} x^{(j)}\big(n^{(j)}\big)^2 \ \le\ 0 ,
\]
так как $x^{(j)} \ge 0$. Равенство требует обнуления всех слагаемых, то есть
\[
  n^{(j)} \ge 0 \ \ \forall j \in N
  \qquad\text{и}\qquad
  n^{(j)}x^{(j)} = 0 \ \ \forall j \in N . \ \blacksquare
\]
Это и есть \textbf{условия Куна--Таккера} ($z = n$: $z \ge 0$, $z\T x = 0$,
$z_B = 0$); при выпуклой $\varphi$ они достаточны для глобального минимума.
Кроме того, $d_N = 0 \Rightarrow d_B = 0$, то есть останов при $d = 0$
корректен.

\paragraph{Выбор шага.}
Равенства сохраняются автоматически; нужна неотрицательность:
\[
  \alpha_{\max} = \min_{i\,:\ d_i < 0}\left(-\frac{x_i}{d_i}\right)
  \qquad (+\infty, \text{ если все } d_i \ge 0) ,
\]
\[
  \alpha_k = \argmin_{0 \le \alpha \le \alpha_{\max}}
    \varphi(x^k + \alpha d^k) .
\]
При $d_i \ge 0$ очевидно $x_i + \alpha d_i \ge 0$; при $d_i < 0$ из
$\alpha \le x_i/(-d_i)$ следует $\alpha(-d_i) \le x_i$, то есть
$x_i + \alpha d_i \ge 0$.

\paragraph{Алгоритм.}
\begin{csteps}
  \item Задать допустимую $x^0 \ge 0$, $Ax^0 = b$, выбрать базис $B$.
  \item Вычислить $\nabla\varphi(x^k)$, приведённый градиент $n_N$, затем
        $d_N$ по правилу Вульфа и $d_B = -B^{-1}Nd_N$.
  \item $d = 0$ — конец (в выпуклом случае точка оптимальна).
  \item Найти $\alpha_{\max}$ и $\alpha_k$; положить
        $x^{k+1} = x^k + \alpha_k d$.
  \item Если базисная переменная обратилась в нуль — заменить её небазисной с
        наибольшей компонентой, пересчитать $B$. К~шагу~2.
\end{csteps}

\paragraph{Теоремы о сходимости.}
\textit{Последовательность $\varphi(x^k)$ убывает; любая предельная точка
$\{x^k\}$ — оптимальное решение.}

\textsc{Схема.} $\varphi(x^k)$ убывает и ограничена снизу $\Rightarrow$
сходится. Базисов конечное число, поэтому берём подпоследовательность с
фиксированным базисом; при фиксированном базисе направление $d(x)$ непрерывно
по $x$. Если для предельной точки $\bar{x}$ было бы $d(\bar{x}) \ne 0$, то
нашёлся бы шаг $\bar\alpha > 0$ с
$\varphi(\bar{x}+\bar\alpha d(\bar{x})) < \varphi(\bar{x})$, что противоречит
сходимости $\varphi(x^k)$. Значит $d(\bar{x}) = 0$: точка стационарна, а при
выпуклой $\varphi$ оптимальна. $\blacksquare$

\paragraph{Ограничения.}
Нужна допустимая начальная точка; требуется выбирать и обновлять базис; при
вырождении возможны трудности; плохо обусловленная $B$ даёт неустойчивость;
при $\alpha_k = \alpha_{\max}$ базис меняется; для невыпуклой $\varphi$ —
только локальный минимум.

\seefull{гл.~3, \S14 (численный пример там же)}

\section{Билет 18. Выпуклый симплексный метод Зангвилла}

\ticketq{Выпуклый симплексный метод Зангвилла.}

\paragraph{Постановка.}
\[
  \varphi(x) \to \min , \qquad Ax = b , \qquad x \ge 0 ,
\]
$\varphi$ выпукла и дифференцируема, $A$ размера $m \times n$,
$\rank A = m$. Разбиение на базис и небазис: $x = (x_B, x_N)$, $A = (B, N)$,
$B$ невырождена, $x_B > 0$ (невырожденный план).

\textbf{Идея.} Обобщение симплекс-метода ЛП на выпуклую нелинейную цель: как и
там, на каждой итерации меняется ровно \emph{одна небазисная} переменная, а
базисные пересчитываются из $Ax = b$.

\paragraph{Приведённый градиент.}
\[
  r_j = \frac{\partial\varphi}{\partial x_j}
      - \nabla_B\T\varphi\; B^{-1}A_j , \qquad j \in N .
\]

\paragraph{Критерий оптимальности (Куна--Таккера).}
\[
  r_j \ge 0 \ \text{при } x_j = 0 ,
  \qquad
  r_j = 0 \ \text{при } x_j > 0 ,
  \qquad j \in N .
\]
Для выпуклой $\varphi$ эти условия достаточны для глобального минимума.

\paragraph{Выбор переменной.}
\[
  u = \min_{j \in N} r_j = r_s ,
  \qquad
  v = \max_{j \in N} x_j r_j = x_t r_t .
\]
\textbf{Останов:} $u \ge 0$ и $v = 0$ — это и есть условия Куна--Таккера
($r_j \ge 0$ всюду и $r_jx_j = 0$).

Иначе: если $\abs{u} \ge v$ — \emph{увеличиваем} $x_s$ ($\sigma = +1$, $q = s$);
если $\abs{u} < v$ — \emph{уменьшаем} $x_t$ ($\sigma = -1$, $q = t$).

\paragraph{Направление.}
\[
  d_q = \sigma , \qquad d_j = 0 \ \ (j \in N,\ j \ne q) ,
  \qquad d_B = -\sigma B^{-1}A_q .
\]

\textbf{Допустимость по равенствам:}
$Ad = Bd_B + A_q d_q = -\sigma A_q + \sigma A_q = 0$ $\Rightarrow$
$A(x + \alpha d) = b$.

\textbf{Спуск.} Так как меняется одна небазисная компонента,
$\nabla\T\varphi\,d = r_q d_q$, и разбор двух случаев:
\begin{itemize}[nosep,left=0pt]
  \item $\abs{u} \ge v$, $\sigma = +1$: тогда $r_q = u < 0$ (иначе был бы
        останов), значит $\nabla\T\varphi\,d = r_q < 0$; увеличение $x_q \ge 0$
        допустимо всегда;
  \item $\abs{u} < v$, $\sigma = -1$: тогда $r_q > 0$ и $x_q > 0$ (иначе
        $v = 0$), значит $\nabla\T\varphi\,d = -r_q < 0$; уменьшать $x_q$ можно
        до нуля.
\end{itemize}
Базисные компоненты строго положительны, поэтому при малом $\alpha$ точка
остаётся допустимой.

\paragraph{Выбор шага.}
\[
  \alpha_{\max} = \min\Big\{ -\frac{x_i}{d_i} \ :\ d_i < 0 \Big\}
  \quad(\text{при } \sigma = -1 \text{ учитывается и } x_q) ,
\]
\[
  \alpha_k = \argmin_{0 \le \alpha \le \alpha_{\max}} \varphi(x_k + \alpha d) .
\]

\paragraph{Алгоритм.}
\begin{csteps}
  \item Задать допустимый базисный план $x^0$ с базисом $B$.
  \item Вычислить $r_j$, $j \in N$; найти $u = r_s$ и $v = x_t r_t$.
  \item $u \ge 0$ и $v = 0$ — конец, $x^k$ оптимальна.
  \item $\abs{u} \ge v$: $\sigma = +1$, $q = s$; иначе $\sigma = -1$, $q = t$.
  \item Построить $d$: $d_q = \sigma$, прочие небазисные нули,
        $d_B = -\sigma B^{-1}A_q$.
  \item Найти $\alpha_{\max}$ и $\alpha_k$; положить
        $x^{k+1} = x^k + \alpha_k d$.
  \item Если на границе базисная переменная обратилась в нуль — она выходит из
        базиса, $x_q$ входит; пересчитать $B$. К~шагу~2.
\end{csteps}

\paragraph{Теорема о сходимости.}
\textit{$\varphi(x^k)$ убывает; предельные точки $\{x^k\}$ — точки
Куна--Таккера, а при выпуклой $\varphi$ — решения задачи.}

\textsc{Схема.} $\varphi(x^k)$ убывает и ограничена снизу $\Rightarrow$
сходится. Число базисов конечно, поэтому можно выделить подпоследовательность
с фиксированным базисом; при фиксированном базисе направление $d(x)$
непрерывно по $x$. Если для предельной точки $\bar{x}$ было бы
$d(\bar{x}) \ne 0$, нашёлся бы шаг $\bar\alpha > 0$ с
$\varphi(\bar{x} + \bar\alpha d(\bar{x})) < \varphi(\bar{x})$, что
противоречит сходимости $\varphi(x^k)$. Значит $d(\bar{x}) = 0$, то есть
выполнены условия Куна--Таккера. $\blacksquare$

\paragraph{Отличие от метода условного градиента.}
Оба метода работают на многограннике $Ax = b$, $x \ge 0$ и используют
приведённый градиент, но вспомогательная задача разная. В методе
Франка--Вульфа (билет~12) решается \emph{задача ЛП на всём} $S$, и $y_k$ —
вершина многогранника, то есть за шаг меняются сразу все координаты. У
Зангвилла задачи ЛП нет вовсе: по правилу $u$/$v$ выбирается \emph{одна}
небазисная переменная, и шаг делается вдоль одного ребра — ровно как в
симплекс-методе ЛП.

\seefull{гл.~3, \S15}

\section{Билет 19. Квадратичное программирование в методе SVM}

\ticketq{Алгоритмы решения задачи квадратичного программирования в методе SVM
(Support Vector Machine — машина опорных векторов).}

\paragraph{Задача классификации.}
Обучающая выборка $(x_i, y_i)$, $i = \overline{1,l}$, $x_i \in \R^n$,
$y_i \in \{-1, 1\}$. Строится разделяющая гиперплоскость $w\T x + b = 0$;
классификация новой точки — $\sgn(w\T x + b)$. Разделение правильное, если
\[
  y_i\,(w\T x_i + b) \ \ge\ 1 , \qquad i = \overline{1,l} .
\]

\textbf{Зазор.} Расстояние от $x_i$ до гиперплоскости равно
$\abs{w\T x_i + b}/\norm{w}$; при выполнении условий выше минимальное
расстояние равно $1/\norm{w}$, а ширина полосы между $w\T x + b = \pm1$ равна
$2/\norm{w}$. Максимизация зазора $\iff$ минимизация $\tfrac12\norm{w}^2$.

\paragraph{Прямая задача.}
Линейно разделимый случай и случай \textbf{мягкого зазора} с переменными
ошибок $\xi_i \ge 0$:
\[
  \tfrac12\norm{w}^2 + C\sum_{i=1}^{l}\xi_i \ \to\ \min_{w,b,\xi} ,
  \qquad
  y_i(w\T x_i + b) \ge 1 - \xi_i , \qquad \xi_i \ge 0 ,
\]
$C > 0$ — параметр штрафа за ошибки. Цель квадратичная, ограничения линейные
$\Rightarrow$ это \textbf{задача квадратичного программирования}.

\paragraph{Вывод двойственной задачи.}
Множители $\alpha_i \ge 0$ для $1 - \xi_i - y_i(w\T x_i + b) \le 0$ и
$\mu_i \ge 0$ для $-\xi_i \le 0$:
\[
  L = \tfrac12\norm{w}^2 + C\sum_i \xi_i
    + \sum_i \alpha_i\big(1 - \xi_i - y_i(w\T x_i + b)\big)
    - \sum_i \mu_i\xi_i .
\]
Условия стационарности:
\[
  \nabla_w L = w - \sum_i \alpha_i y_i x_i = 0
  \ \Rightarrow\ w = \sum_i \alpha_i y_i x_i ;
\]
\[
  \frac{\partial L}{\partial b} = -\sum_i \alpha_i y_i = 0
  \ \Rightarrow\ \sum_i \alpha_i y_i = 0 ;
  \qquad
  \frac{\partial L}{\partial \xi_i} = C - \alpha_i - \mu_i = 0
  \ \Rightarrow\ 0 \le \alpha_i \le C .
\]
Подстановка $w$ в $L$ даёт двойственную задачу (с ядром
$K(x_i,x_j) = \varphi\T(x_i)\varphi(x_j)$):
\[
  W(\alpha) = \sum_{i} \alpha_i
   - \tfrac12\sum_{i}\sum_{j}\alpha_i\alpha_j y_i y_j K(x_i,x_j)
   \ \to\ \max_\alpha ,
\]
\[
  \sum_i \alpha_i y_i = 0 , \qquad 0 \le \alpha_i \le C .
\]
Решающая функция $F(x) = \sum_i \alpha_i y_i K(x_i, x) + b$, классификация
$\sgn F(x)$. На практике решают именно двойственную: она зависит только от
скалярных произведений (позволяет ядра), имеет простые ограничения, и
большинство $\alpha_i$ равны нулю.

\paragraph{Условия ККТ и опорные векторы.}
\[
  \alpha_i\big(y_i(w\T x_i + b) - 1 + \xi_i\big) = 0 ,
  \qquad (C - \alpha_i)\xi_i = 0 .
\]
Отсюда: $\alpha_i = 0$ — объект не влияет на гиперплоскость;
$0 < \alpha_i < C$ — объект лежит на границе зазора,
$y_i(w\T x_i + b) = 1$; $\alpha_i = C$ — объект нарушает зазор или
классифицирован неверно. Объекты с $\alpha_i > 0$ — \textbf{опорные векторы}.

\paragraph{Теорема (выпуклость двойственной задачи).}
\textit{Двойственная задача SVM — выпуклая задача квадратичного
программирования.}

\textsc{Док-во.} $W(\alpha) = \sum_i\alpha_i - \tfrac12\alpha\T Q\alpha$,
$Q_{ij} = y_iy_jK(x_i,x_j)$. Для любого $z$, положив $u_i = z_iy_i$:
\[
  z\T Q z = \sum_i\sum_j z_iz_jy_iy_jK(x_i,x_j)
          = \sum_i\sum_j u_iu_jK(x_i,x_j) \ \ge\ 0 ,
\]
так как ядровая матрица Грама положительно полуопределена $\Rightarrow$
$Q \succeq 0$ $\Rightarrow$ $-\tfrac12\alpha\T Q\alpha$ вогнута, линейная
часть вогнута (и выпукла) $\Rightarrow$ $W$ вогнута. Максимизация вогнутой
функции при линейных ограничениях — выпуклая задача. $\blacksquare$

\paragraph{Алгоритмы решения.}
\begin{csteps}
  \item \textbf{Метод активного набора.} Часть ограничений ($\alpha_i = 0$ или
        $\alpha_i = C$) считается активной; на итерации решается КП при
        активных ограничениях, проверяются условия ККТ, набор пополняется или
        сокращается. Минусы: хранение матрицы ядра, частая смена набора при
        большом $l$.
  \item \textbf{Метод внутренней точки.} Ограничения $0 < \alpha_i < C$
        заменяются логарифмическим барьером
        $-\mu\sum_i\ln\alpha_i - \mu\sum_i\ln(C - \alpha_i)$, решается
        последовательность гладких задач внутри области. Минусы: большие
        системы линейных уравнений, тяжёл при больших выборках.
  \item \textbf{Методы декомпозиции.} Оптимизируются не все $\alpha$ сразу, а
        небольшое рабочее множество $B_k \subset \{1,\dots,l\}$; остальные
        фиксируются.
\end{csteps}

\paragraph{Алгоритм SMO.}
Частный случай декомпозиции: на итерации меняются ровно \emph{две} переменные
$\alpha_i, \alpha_j$. Две — потому что есть равенство
$\sum_s \alpha_s y_s = 0$: изменение одной переменной его нарушит.

\textbf{Теорема (сохранение равенства).} \textit{Если
$y_i\alpha_i^{\text{new}} + y_j\alpha_j^{\text{new}}
 = y_i\alpha_i^{\text{old}} + y_j\alpha_j^{\text{old}}$, то
$\sum_s \alpha_s y_s = 0$ сохраняется.}

\textsc{Док-во.} Меняются только компоненты $i$ и $j$, поэтому
\[
  \sum_s \alpha_s^{\text{new}}y_s - \sum_s \alpha_s^{\text{old}}y_s
   = y_i(\alpha_i^{\text{new}} - \alpha_i^{\text{old}})
   + y_j(\alpha_j^{\text{new}} - \alpha_j^{\text{old}}) = 0
\]
по условию теоремы $\Rightarrow$ новая сумма равна старой, то есть нулю.
$\blacksquare$

\textbf{Откуда формула обновления.} При фиксированных остальных $\alpha$ связь
$\sum_t\alpha_ty_t = \mathrm{const}$ даёт $\alpha_i = \zeta - y_iy_j\alpha_j$,
то есть $W$ становится функцией одной переменной $\alpha_j$ — параболой:
\[
  \frac{dW}{d\alpha_j} = y_j(E_i - E_j) ,
  \qquad
  \frac{d^2W}{d\alpha_j^2} = -\eta \le 0 ,
  \qquad
  \eta = K_{ii} + K_{jj} - 2K_{ij} ,
\]
где $F(x) = \sum_s\alpha_sy_sK(x_s,x) + b$ и $E_t = F(x_t) - y_t$. Вершина
параболы и есть новое значение.

\textbf{Формулы.}
\[
  \alpha_j^{\text{new}} = \alpha_j^{\text{old}}
    + \frac{y_j(E_i - E_j)}{\eta} \quad (\eta > 0) ,
\]
затем отсечение $L \le \alpha_j^{\text{new}} \le H$:
\[
  y_i \ne y_j :\ L = \max(0,\ \alpha_j - \alpha_i) ,\
                 H = \min(C,\ C + \alpha_j - \alpha_i) ;
\]
\[
  y_i = y_j :\ L = \max(0,\ \alpha_i + \alpha_j - C) ,\
               H = \min(C,\ \alpha_i + \alpha_j) ,
\]
и $\alpha_i^{\text{new}} = \alpha_i^{\text{old}}
 + y_iy_j\big(\alpha_j^{\text{old}} - \alpha_j^{\text{new}}\big)$ — ровно
условие теоремы. Свободный член:
$b^{\text{new}} = b_1$ при $0 < \alpha_i^{\text{new}} < C$,
$b^{\text{new}} = b_2$ при $0 < \alpha_j^{\text{new}} < C$, иначе
$(b_1+b_2)/2$, где
\[
  b_1 = b - E_i - y_i\Delta\alpha_i K(x_i,x_i) - y_j\Delta\alpha_j K(x_i,x_j) ,
\]
\[
  b_2 = b - E_j - y_i\Delta\alpha_i K(x_i,x_j) - y_j\Delta\alpha_j K(x_j,x_j) .
\]

\paragraph{Ограничения.}
Матрица ядра имеет размер $l\times l$ — при больших $l$ её хранение
невозможно; активные методы и методы внутренней точки для таких выборок
слишком дороги; SMO делает дешёвые итерации, но их может потребоваться много,
и качество зависит от выбора пары $(i,j)$; при плохом $C$ модель
переобучается или недообучается; при неудачном ядре данные плохо разделяются.

\seefull{гл.~4, \S1--\S4 (численный пример SMO в \S4)}


\end{document}
